% Rozdzial do wlaczenia do skryptu: % Rozdzial do wlaczenia do skryptu: % Rozdzial do wlaczenia do skryptu: % Rozdzial do wlaczenia do skryptu: \input{geometrie_thurstona_3rozmaitosci}
% Wymagane w preambule: amsmath, amssymb, amsthm, mathtools.
% Srodowiska: definicja, twierdzenie, przyklad, uwaga, wniosek, proof.
% Jezeli w skrypcie stosowana jest inna definicja srodowisk,
% nalezy dostosowac nazwy w preambule, nie w tresci rozdzialu.

\chapter{Geometrie Thurstona i struktury geometryczne na trójrozmaitościach}
\label{chap:geometrie-thurstona}

W poprzednich rozdziałach rozważaliśmy rozmaitości Riemanna, ich połączenia,
geodezyjne i krzywiznę oraz podstawowe geometrie przestrzeni o stałej krzywiźnie.
Tutaj przejdziemy od geometrii pojedynczej rozmaitości do pytania o to,
\emph{jakie geometrie jednorodne mogą występować lokalnie na trójrozmaitościach}.
Z jednej strony prowadzi to do klasyfikacji ośmiu geometrii modelowych
Thurstona, z drugiej -- do związków między metryką a topologią rozmaitości
trójwymiarowych. W szczególności zobaczymy, dlaczego do opisu wszystkich
zwartych trójrozmaitości nie wystarczają geometrie sferyczna, euklidesowa
i hiperboliczna.

\section{Struktury \texorpdfstring{$(G,X)$}{(G,X)} i geometrie modelowe}

Ideą programu erlangeńskiego Kleina jest badanie przestrzeni za pomocą
grupy przekształceń zachowujących interesującą nas strukturę. Geometria
Riemanna dostarcza szczególnego przypadku: jeśli $(X,g)$ jest rozmaitością
Riemanna, to grupa $\operatorname{Isom}(X,g)$ składa się z jej izometrii.

\begin{definicja}[Struktura $(G,X)$]
Niech $X$ będzie spójną gładką rozmaitością, a $G$ grupą Liego działającą na
$X$ przez dyfeomorfizmy. \emph{Strukturą $(G,X)$} na rozmaitości $M$
nazywamy atlas map $\{(U_\alpha,\varphi_\alpha)\}$,
gdzie $\varphi_\alpha\colon U_\alpha\to V_\alpha\subset X$ jest
dyfeomorfizmem na zbiór otwarty, a dla każdej spójnej składowej
$C\subset U_\alpha\cap U_\beta$ istnieje $g_{\alpha\beta,C}\in G$ takie, że
$\varphi_\alpha\circ\varphi_\beta^{-1}$ na $\varphi_\beta(C)$
jest obcięciem działania $g_{\alpha\beta,C}$.
\end{definicja}

Warunek dotyczący składowych przecięcia jest istotny: jedno przecięcie
$U_\alpha\cap U_\beta$ może być niespójne, a na różnych jego składowych
mogą występować różne elementy grupy $G$. Zwykły atlas gładki odpowiada
pseudogrupie wszystkich lokalnych dyfeomorfizmów, podczas gdy struktura
$(G,X)$ dopuszcza tylko lokalne przekształcenia pochodzące z jednej grupy.

\begin{przyklad}[Torus płaski]
Niech $X=\mathbb R^3$, $G=\operatorname{Isom}(\mathbb R^3)$ i
$\Gamma=\mathbb Z^3$ działająca przez translacje. Iloraz
$\mathbb T^3=\mathbb R^3/\mathbb Z^3$ otrzymuje atlas $(G,X)$, ponieważ
lokalne odwrotności nakrycia różnią się o translacje należące do $G$.
Wszystkie przestrzenie styczne są euklidesowe, chociaż $\mathbb T^3$
nie jest globalnie izometryczny z $\mathbb R^3$.
\end{przyklad}

\begin{definicja}[Geometria modelowa Thurstona]
\label{def:model-thurston}
Przez \emph{geometrię modelową} rozumiemy parę $(G,X)$, dla której:
\begin{enumerate}
\item $X$ jest spójna i jednospójna;
\item $G$ działa na $X$ tranzytywnie, a stabilizator punktu
      $G_x=\{a\in G:a\cdot x=x\}$ jest zwarty;
\item $G$ jest maksymalną grupą dyfeomorfizmów $X$ spełniającą warunek
      zwartości stabilizatorów;
\item istnieje zwarta rozmaitość z pełną strukturą $(G,X)$.
\end{enumerate}
Dwie geometrie utożsamiamy, gdy ich przestrzenie i grupy działań są
izomorficzne w naturalnym sensie. Przymiotnik \emph{pełna} oznacza, że
struktura pochodzi z ilorazu $X/\Gamma$ dla odpowiednio działającej
grupy dyskretnej $\Gamma\subset G$.
\end{definicja}

Dwa ostatnie warunki nie są kosmetyczne. Gdyby zrezygnować z maksymalności,
tę samą geometrię euklidesową można by sztucznie traktować jako wiele
różnych geometrii, wybierając mniejsze grupy ruchów. Wymóg zwartego
ilorazu eliminuje z kolei liczne jednorodne przestrzenie niemające
zwartych realizacji globalnych.

\begin{twierdzenie}[Metryka niezmiennicza]
\label{thm:metryka-niezmiennicza}
Jeżeli grupa Liego $G$ działa tranzytywnie na $X$, a stabilizator $G_o$
pewnego punktu $o\in X$ jest zwarty, to na $X$ istnieje metryka
Riemanna niezmiennicza względem $G$.
\end{twierdzenie}
\begin{proof}
Wybierzmy dowolny iloczyn skalarny $h$ na $T_oX$. Niech $\mu$ będzie
unormowaną miarą Haara na zwartej grupie $G_o$. Definiujemy
\[
 h_o^{G}(u,v)=\int_{G_o}h(da_o u,da_o v)\,d\mu(a).
\]
Otrzymana forma jest dodatnio określona i niezmiennicza względem
reprezentacji izotropii. Jeśli $x=b\cdot o$, przenosimy ją wzorem
$h_x^G(db_o u,db_o v)=h_o^G(u,v)$. Definicja nie zależy od wyboru $b$:
dwa takie elementy różnią się o element stabilizatora $G_o$, a forma
$h_o^G$ została względem niego uśredniona. Gładkość metryki wynika z
gładkości działania i lokalnego wyboru przekrojów odwzorowania
$G\to G/G_o\simeq X$.
\end{proof}

\begin{uwaga}[Nakrycie uniwersalne i grupa decka]
Jeżeli $X$ jest jednospójna, $\Gamma\subset\operatorname{Isom}(X,g)$
działa swobodnie, właściwie nieciągle i współzwarto, to
$M=\Gamma\backslash X$ jest zwartą rozmaitością Riemanna, a
$X\to M$ jest jej nakryciem uniwersalnym. W szczególności
$\pi_1(M)\simeq\Gamma$. Odwrotnie, każda \emph{pełna} struktura
$(G,X)$ na zwartej rozmaitości jest realizowana w ten sposób.
Nie każda struktura $(G,X)$ jest jednak z definicji pełna.
\end{uwaga}

Przydatne pojęcie pośrednie daje \emph{odwzorowanie rozwijające}
$\operatorname{dev}\colon\widetilde M\to X$. Powstaje ono przez analityczne
kontynuowanie kolejnych lokalnych map modelowych na nakryciu
uniwersalnym i spełnia
$\operatorname{dev}(\gamma p)=\rho(\gamma)\operatorname{dev}(p)$,
gdzie $\rho\colon\pi_1(M)\to G$ jest \emph{holonomią}. Jeżeli struktura
jest pełna, odwzorowanie rozwijające identyfikuje $\widetilde M$ z $X$.

\section{Dlaczego w wymiarze trzecim pojawia się osiem geometrii?}

\begin{twierdzenie}[Klasyfikacja geometrii modelowych Thurstona]
\label{thm:osiem-geometrii}
W wymiarze trzecim, z dokładnością do izomorfizmu geometrii modelowych,
istnieje dokładnie osiem możliwości:
\[
\mathbb S^3,\qquad \mathbb R^3,\qquad \mathbb H^3,\qquad
\mathbb S^2\!\times\!\mathbb R,\qquad
\mathbb H^2\!\times\!\mathbb R,\qquad
\mathrm{Nil},\qquad\widetilde{\mathrm{SL}}_2(\mathbb R),\qquad
\mathrm{Sol}.
\]
Dla trzech pierwszych geometrii spójna izotropia jest trójwymiarowa;
dla następnych czterech -- jednowymiarowa; dla geometrii Sol --
zerowymiarowa.
\end{twierdzenie}

Podamy mechanizm klasyfikacji, oddzielając go od pełnego dowodu, który
wymaga dodatkowej analizy jednorodnych grup Liego i ich sieci.
Na mocy twierdzenia~\ref{thm:metryka-niezmiennicza} możemy wybrać
metrykę niezmienniczą na $X$. Izotropia działa wiernie na $T_oX$
(izometria ustalająca punkt i mająca identyczną różniczkę jest
identycznością), zatem jej składowa spójna jest zwartą spójną
podgrupą $\mathrm{SO}(3)$. Takie podgrupy są, z dokładnością do
sprzężenia, równe $\mathrm{SO}(3)$, $\mathrm{SO}(2)$ albo grupie
trywialnej.

Gdy izotropia ma wymiar trzy, działa tranzytywnie na kierunkach
w każdej przestrzeni stycznej. Krzywizna przekrojowa zależy więc
tylko od punktu, a jednorodność czyni ją stałą na całym $X$.
Otrzymujemy trzy przestrzenie stałej krzywizny.

Gdy izotropia ma wymiar jeden, zachowuje pewien wyróżniony kierunek
pola prostych. Na odpowiednim nakryciu pole to można zorientować i
otrzymać niezmiennicze pole wektorowe, którego linie całkowe tworzą
jednowymiarowe włóknowanie. Geometria poprzeczna jest jedną z trzech
geometrii powierzchni jednorodnych; krzywizna połączenia tego
włóknowania odróżnia geometrie produktowe od Nil i
$\widetilde{\mathrm{SL}}_2(\mathbb R)$. Przy zerowymiarowej izotropii
analizuje się trójwymiarowe grupy Liego wyposażone w metryki
lewostronnie niezmiennicze. Po uwzględnieniu maksymalności grupy
izometrii i istnienia zwartego ilorazu jedyną nową geometrią jest Sol.
To jest \emph{szkic dowodu}, nie pełna klasyfikacja możliwych algebr
Liego i ich sieci.

\begin{center}
\begin{tabular}{|l|c|l|}
\hline
Geometria & $\dim G_o^0$ & Typ geometrii poprzecznej lub krzywizna \\
\hline
$\mathbb S^3$ & $3$ & $K\equiv 1$ \\
$\mathbb R^3$ & $3$ & $K\equiv 0$ \\
$\mathbb H^3$ & $3$ & $K\equiv -1$ \\
$\mathbb S^2\times\mathbb R$ & $1$ & produkt, baza $\mathbb S^2$ \\
$\mathbb H^2\times\mathbb R$ & $1$ & produkt, baza $\mathbb H^2$ \\
Nil & $1$ & baza $\mathbb R^2$, skręcone włókna \\
$\widetilde{\mathrm{SL}}_2(\mathbb R)$ & $1$ & baza $\mathbb H^2$, skręcone włókna \\
Sol & $0$ & włóknowanie nad $\mathbb R$ \\
\hline
\end{tabular}
\end{center}

\section{Trzy geometrie o stałej krzywiźnie}

Przypomnijmy jedynie elementy potrzebne do dalszego porównania. Po
odpowiednim przeskalowaniu metryk geometrie te mają krzywiznę
przekrojową odpowiednio $\kappa=1,0,-1$. W wymiarze trzy ich
tensor Ricciego i krzywizna skalarna spełniają
$\operatorname{Ric}=2\kappa g$ oraz $\operatorname{Scal}=6\kappa$.

\begin{przyklad}[Modele i zwarte ilorazy]
\begin{enumerate}
\item $\mathbb S^3=\{z\in\mathbb C^2:|z_1|^2+|z_2|^2=1\}$,
      z metryką indukowaną z $\mathbb R^4$. Przykładami zwartych
      ilorazów są przestrzenie soczewkowe
      $L(p,q)=\mathbb S^3/\mathbb Z_p$, gdzie generator działa przez
      $(z_1,z_2)\mapsto(e^{2\pi i/p}z_1,e^{2\pi iq/p}z_2)$,
      a $\gcd(p,q)=1$.
\item $\mathbb R^3$, z metryką $g=dx^2+dy^2+dz^2$;
      ilorazem zwartym jest $\mathbb T^3$.
\item $\mathbb H^3=\{(x,y,z)\in\mathbb R^3:z>0\}$,
      z metryką $g=z^{-2}(dx^2+dy^2+dz^2)$. Istnieją zwarte ilorazy
      $\Gamma\backslash\mathbb H^3$, gdzie $\Gamma$ jest bezskrętną
      współzwartą kratą w $\operatorname{Isom}(\mathbb H^3)$.
\end{enumerate}
\end{przyklad}

\section{Dwie geometrie produktowe}

Niech $(\Sigma^2,g_\Sigma)$ będzie powierzchnią o stałej krzywiźnie
Gaussa $\kappa\in\{1,-1\}$. Na $X=\Sigma\times\mathbb R$ bierzemy
metrykę produktową $g=g_\Sigma+dt^2$. Ponieważ połączenie
Levi-Civity produktu zachowuje oba czynniki, dla ortonormalnych
wektorów $e_1,e_2$ stycznych do $\Sigma$ oraz $e_3=\partial_t$
mamy $K(e_1,e_2)=\kappa$ i $K(e_1,e_3)=K(e_2,e_3)=0$.
W konsekwencji
\[
\operatorname{Ric}(e_1,e_1)=\operatorname{Ric}(e_2,e_2)=\kappa,
\qquad \operatorname{Ric}(e_3,e_3)=0,
\qquad \operatorname{Scal}=2\kappa.
\]
Są to geometrie jednorodne, lecz \emph{nie} mają stałej krzywizny
przekrojowej w wymiarze trzy.

\begin{przyklad}[Zwarte realizacje]
Dla $\mathbb S^2\times\mathbb R$ otrzymujemy
$\mathbb S^2\times\mathbb S^1$ przez ilorazowanie drugiego czynnika
translacjami całkowitymi. Dla $\mathbb H^2\times\mathbb R$ wybieramy
zwartą powierzchnię hiperboliczną $\Sigma_g=\Gamma\backslash\mathbb H^2$
rodzaju $g\ge2$; wtedy $\Sigma_g\times\mathbb S^1$ jest zwartym
ilorazem $\mathbb H^2\times\mathbb R$.
\end{przyklad}

\section{Geometria Nil: grupa Heisenberga}
\label{sec:thurston-nil}

Pierwsza istotnie nowa geometria pojawia się wtedy, gdy płaszczyznę
$\mathbb R^2$ zastąpimy rodziną prostych, których połączenie ma
niezerową krzywiznę. Rozważmy $\mathbb R^3$ z działaniem grupowym
\[
(x,y,z)\cdot(x',y',z')=
(x+x',\,y+y',\,z+z'+xy').
\]
Jest to trójwymiarowa \emph{grupa Heisenberga} $H_3$. Można ją
utożsamić z grupą macierzy
$\left(\begin{smallmatrix}1&x&z\\0&1&y\\0&0&1\end{smallmatrix}\right)$.

Jednoformy $\theta^1=dx$, $\theta^2=dy$ i
$\theta^3=dz-x\,dy$ są lewostronnie niezmiennicze.
Ich formami dualnymi są pola
$E_1=\partial_x$, $E_2=\partial_y+x\partial_z$ oraz
$E_3=\partial_z$. Bezpośrednio
$[E_1,E_2]=E_3$, zaś pozostałe nawiasy są zerowe.
Nadajmy $H_3$ metrykę
\begin{equation}\label{eq:nil-metric}
 g_{\mathrm{Nil}}=(\theta^1)^2+(\theta^2)^2+(\theta^3)^2
 =dx^2+dy^2+(dz-x\,dy)^2.
\end{equation}
Zapis $\theta^2$ w powyższej formule oznacza iloczyn tensorowy
$\theta\otimes\theta$, a nie iloczyn zewnętrzny. Pola $E_1,E_2,E_3$
tworzą globalną ortonormalną ramę.

\begin{przyklad}[Pozioma krzywa i holonomia]
Rzut $\pi\colon H_3\to\mathbb R^2$, $\pi(x,y,z)=(x,y)$,
ma włókna równoległe do osi $z$. Rozkład poziomy określa
$\theta^3=0$, czyli dla krzywej $\gamma(s)=(x(s),y(s),z(s))$
warunek poziomości brzmi $z'(s)=x(s)y'(s)$. Po przejściu
wzdłuż zamkniętej pętli $c$ w płaszczyźnie podstawa zostaje
niezmieniona, ale współrzędna $z$ zmienia się o
\[
\Delta z=\oint_c x\,dy=\iint_D dx\wedge dy,
\]
gdzie $D$ jest zorientowanym obszarem ograniczonym przez $c$.
Oznacza to, że przesunięcie wzdłuż włókna mierzy pole zorientowane
ograniczone przez pętlę. Ponieważ $d\theta^3=-dx\wedge dy$,
rozkład poziomy nie jest całkowalny.
\end{przyklad}

\begin{twierdzenie}[Krzywizny geometrii Nil]
Dla metryki~\eqref{eq:nil-metric} krzywizny przekrojowe płaszczyzn
wyznaczonych przez ramę $E_1,E_2,E_3$ wynoszą
$K_{12}=-3/4$ i $K_{13}=K_{23}=1/4$.
Stąd $\operatorname{Ric}=\operatorname{diag}(-1/2,-1/2,1/2)$
w tej ramie i $\operatorname{Scal}=-1/2$.
\end{twierdzenie}
\begin{proof}
Dla ortonormalnych pól lewostronnie niezmienniczych wzór Koszula
przybiera postać
\[
2\langle\nabla_XY,Z\rangle
=\langle[X,Y],Z\rangle-\langle[Y,Z],X\rangle
+\langle[Z,X],Y\rangle.
\]
Z jedynego niezerowego nawiasu $[E_1,E_2]=E_3$ dostajemy
\[
\begin{aligned}
&\nabla_{E_1}E_2=\tfrac12E_3,\qquad
 \nabla_{E_2}E_1=-\tfrac12E_3,\\
&\nabla_{E_1}E_3=\nabla_{E_3}E_1=-\tfrac12E_2,\\
&\nabla_{E_2}E_3=\nabla_{E_3}E_2=\tfrac12E_1,
\qquad \nabla_{E_i}E_i=0.
\end{aligned}
\]
Przy konwencji $R(X,Y)Z=\nabla_X\nabla_YZ-\nabla_Y\nabla_XZ-
\nabla_{[X,Y]}Z$ otrzymujemy na przykład
$R(E_1,E_2)E_2=-\tfrac34E_1$ oraz
$R(E_1,E_3)E_3=\tfrac14E_1$. Trzeci rachunek jest analogiczny.
Tensor Ricciego jest sumą krzywizn przekrojowych po ortonormalnych
kierunkach; daje to podane wartości.
\end{proof}

\begin{przyklad}[Zwarta nilrozmaitość]
Podgrupa $\Gamma=H_3(\mathbb Z)$ macierzy Heisenberga o całkowitych
współrzędnych jest dyskretna i współzwarta. Iloraz
$\Gamma\backslash H_3(\mathbb R)$ jest zwartą trójrozmaitością
z geometrią Nil. Stanowi nietrywialną wiązkę okręgów nad
$\mathbb T^2$: przechodząc wokół podstawowego równoległoboku w bazie,
włókno doznaje przesunięcia wynikającego z niezerowej krzywizny
połączenia. Nie jest to torus płaski $\mathbb T^3$.
\end{przyklad}

\section{Geometria \texorpdfstring{$\widetilde{\mathrm{SL}}_2(\mathbb R)$}{SL2-tylda}}
\label{sec:thurston-sl2}

Geometria ta jest naturalnie związana z jednostkową wiązką styczną
płaszczyzny hiperbolicznej. Grupa izometrii zachowujących orientację
$\mathbb H^2$ jest izomorficzna z
$\operatorname{PSL}_2(\mathbb R)$; działa ona swobodnie i tranzytywnie
na wektorach stycznych długości jeden. Stąd
$T^1\mathbb H^2\simeq\operatorname{PSL}_2(\mathbb R)$ jako rozmaitości.
Nakryciem uniwersalnym tej przestrzeni jest
$\widetilde{\operatorname{SL}}_2(\mathbb R)$.

W modelu górnej półpłaszczyzny $\mathbb H^2=\{(x,y):y>0\}$
można, po identyfikacji nakrycia uniwersalnego z
$\mathbb H^2\times\mathbb R$, wybrać metrykę
\begin{equation}\label{eq:sl2-metric}
 g_{\widetilde{\mathrm{SL}}_2}
 =\frac{dx^2+dy^2}{y^2}
 +\left(dt+\frac{dx}{y}\right)^2.
\end{equation}
Rzut na $(x,y)$ jest submersją Riemannowską o jednowymiarowych,
geodezyjnych włóknach. Forma połączenia
$\alpha=dt+dx/y$ ma krzywiznę
$d\alpha=dx\wedge dy/y^2$, a więc jest ona (z uwzględnieniem
orientacji) równa formie pola powierzchni hiperbolicznej.
Nie jest to metryka produktowa.

\begin{uwaga}[Krzywizny]
Dla normalizacji~\eqref{eq:sl2-metric} wzory O'Neilla dla
submersji Riemannowskiej dają krzywiznę poziomej płaszczyzny
$-1-3/4=-7/4$ oraz krzywizny płaszczyzn zawierających kierunek
włókna równe $1/4$. W szczególności w poziomej ortonormalnej
bazie $E_1,E_2$ i pionowym $E_3$ tensor Ricciego ma wartości
$(-3/2,-3/2,1/2)$. Konkretne wartości zależą od normalizacji
współczynnika przy $\alpha^2$; znak niezerowej krzywizny
połączenia jest ważniejszy od tych liczb.
\end{uwaga}

\begin{przyklad}[Wiązka jednostkowych wektorów nad powierzchnią]
Niech $\Sigma_g=\Gamma\backslash\mathbb H^2$ będzie zamkniętą,
zorientowaną powierzchnią hiperboliczną rodzaju $g\ge2$.
Wówczas $T^1\Sigma_g$ jest zwartą trójrozmaitością, która otrzymuje
geometrię $\widetilde{\operatorname{SL}}_2(\mathbb R)$.
Wiązka $\mathbb S^1\to T^1\Sigma_g\to\Sigma_g$ ma liczbę
Eulera $\chi(\Sigma_g)=2-2g\neq0$. Wynika to z faktu, że
$T^1\Sigma_g$ jest wiązką okręgów stowarzyszoną ze zorientowaną
wiązką styczną $T\Sigma_g$, której klasa Eulera paruje z klasą
fundamentalną do charakterystyki Eulera powierzchni.
\end{przyklad}

Powyższy przykład należy porównać z
$\Sigma_g\times\mathbb S^1$ o geometrii
$\mathbb H^2\times\mathbb R$. Obie rozmaitości mają wiązkę okręgów
nad tą samą powierzchnią, ale w pierwszej liczba Eulera wiązki
jest niezerowa, zaś w drugim przypadku zerowa. Topologia
wiązki wykrywa więc różnicę między dwiema geometriami.

\section{Geometria Sol i odwzorowania sklejenia torusa}
\label{sec:thurston-sol}

Grupę Sol można zrealizować jako $\mathbb R^2\rtimes\mathbb R$,
gdzie $t\in\mathbb R$ działa na $\mathbb R^2$ macierzą
$\operatorname{diag}(e^t,e^{-t})$. Prawo grupowe jest dane wzorem
\[
(x,y,t)\cdot(x',y',t')
=(x+e^t x',\,y+e^{-t}y',\,t+t').
\]
W odróżnieniu od Nil, ta grupa jest rozwiązalna, ale nie nilpotentna.
Jednoformy $\theta^1=e^{-t}dx$, $\theta^2=e^t dy$ i
$\theta^3=dt$ są lewostronnie niezmiennicze. Wybieramy metrykę
\begin{equation}\label{eq:sol-metric}
 g_{\mathrm{Sol}}=e^{-2t}dx^2+e^{2t}dy^2+dt^2.
\end{equation}
Jej globalna ortonormalna rama to
$E_1=e^t\partial_x$, $E_2=e^{-t}\partial_y$,
$E_3=\partial_t$; spełnia ona
$[E_3,E_1]=E_1$, $[E_3,E_2]=-E_2$ oraz $[E_1,E_2]=0$.

\begin{twierdzenie}[Krzywizny geometrii Sol]
Dla metryki~\eqref{eq:sol-metric} zachodzi
$K_{12}=1$ oraz $K_{13}=K_{23}=-1$.
W szczególności
$\operatorname{Ric}=\operatorname{diag}(0,0,-2)$
w ramie $E_1,E_2,E_3$ i $\operatorname{Scal}=-2$.
\end{twierdzenie}
\begin{proof}
Ze wzoru Koszula otrzymujemy
$\nabla_{E_1}E_1=E_3$,
$\nabla_{E_1}E_3=-E_1$,
$\nabla_{E_2}E_2=-E_3$,
$\nabla_{E_2}E_3=E_2$;
pozostałe pochodne kowariantne pól ramy są zerowe.
Stąd $R(E_1,E_2)E_2=E_1$ oraz
$R(E_1,E_3)E_3=-E_1$ i $R(E_2,E_3)E_3=-E_2$.
Pozostałe wzory wynikają z definicji krzywizny Ricciego.
\end{proof}

\begin{przyklad}[Zwarta geometria Sol jako torus odwzorowania]
\label{ex:sol-mapping-torus}
Niech
$A=\left(\begin{smallmatrix}2&1\\1&1\end{smallmatrix}\right)
\in\mathrm{SL}_2(\mathbb Z)$. Jej wartości własne są rzeczywiste,
dodatnie i równe
$\lambda=(3+\sqrt5)/2>1$ oraz $\lambda^{-1}$.
Rozważmy \emph{torus odwzorowania}
\[
M_A=\frac{\mathbb T^2\times[0,1]}
{(u,1)\sim(Au,0)}.
\]
Jest to zwarta trójrozmaitość będąca wiązką torusów nad
$\mathbb S^1$. Aby zobaczyć geometrię Sol, wybierzmy liniowe
współrzędne na $\mathbb R^2$ diagonalizujące $A$.
W tych współrzędnych działanie $A$ ma postać
$(x,y)\mapsto(\lambda x,\lambda^{-1}y)$, natomiast
$\mathbb Z^2$ przechodzi w pewną kratę $\Lambda\subset\mathbb R^2$,
niezmienniczą względem tego przekształcenia.
Grupa generowana przez translacje o wektory $\Lambda$ oraz
\[
 (x,y,t)\longmapsto
 (\lambda x,\lambda^{-1}y,t+\log\lambda)
\]
działa na $(\mathrm{Sol},g_{\mathrm{Sol}})$ izometrycznie,
swobodnie i współzwarto. Otrzymany iloraz jest torusem
odwzorowania z monodromią $A$ (ewentualnie $A^{-1}$,
zależnie od konwencji sklejania końców przedziału).
Zatem $M_A$ ma geometrię Sol.
\end{przyklad}

Widać tu bezpośrednio związek algebry z geometrią: hiperboliczne
odwzorowanie torusa ma dwa kierunki własne, z których jeden ulega
rozciąganiu, a drugi ściskaniu. Metryka Sol dokładnie kompensuje te
zmiany przy przesunięciu w kierunku trzeciej współrzędnej.

\section{Wiązki Seiferta a geometrie modelowe}

\begin{definicja}[Wiązka Seiferta]
W uproszczeniu \emph{rozmaitością włóknistą Seiferta} nazywamy
trójrozmaitość rozłożoną na okręgi tak, że każde włókno ma
otoczenie modelowane na odpowiednim włóknistym torusie pełnym.
Włókna wyjątkowe dopuszczają skończone lokalne stabilizatory;
przestrzeń ilorazowa jest wówczas dwuwymiarowym orbifoldem.
Gdy włókien wyjątkowych brak, mamy zwykłą wiązkę okręgów nad
powierzchnią.
\end{definicja}

Przykłady $\mathbb S^2\times\mathbb S^1$,
$\Sigma_g\times\mathbb S^1$, zwarta nilrozmaitość oraz
$T^1\Sigma_g$ pokazują wspólny wzorzec: większość geometrii
Thurstona z jednowymiarową izotropią występuje naturalnie na
rozmaitościach z włóknowaniem przez okręgi.
Odpowiednia geometria zależy od krzywizny orbifoldu bazy
oraz liczby Eulera włóknowania. Dla zorientowanej wiązki
okręgów nad zamkniętą zorientowaną powierzchnią sytuację
ilustruje poniższe zestawienie; w ogólności trzeba uwzględnić
również włókna wyjątkowe.

\begin{center}
\begin{tabular}{|c|c|l|}
\hline
Krzywizna bazy & Liczba Eulera wiązki & Geometria \\
\hline
$+1$ & $0$ & $\mathbb S^2\times\mathbb R$ \\
$0$ & $0$ & $\mathbb R^3$ \\
$0$ & $\neq0$ & Nil \\
$-1$ & $0$ & $\mathbb H^2\times\mathbb R$ \\
$-1$ & $\neq0$ & $\widetilde{\mathrm{SL}}_2(\mathbb R)$ \\
\hline
\end{tabular}
\end{center}

Tabela nie obejmuje sferycznych przestrzeni z niezerową klasą
Eulera, które również bywają włókniste Seiferta. Pełne kryteria
dla orbifoldów wykorzystują ich orbifoldową charakterystykę
Eulera oraz odpowiednio znormalizowaną liczbę Eulera
włóknowania. W wersji ogólnej dla zamkniętych, orientowalnych
rozmaitości Seiferta przy $\chi^{\mathrm{orb}}>0$ występuje
geometria $\mathbb S^3$ lub $\mathbb S^2\times\mathbb R$;
przy $\chi^{\mathrm{orb}}=0$ geometria Nil lub $\mathbb R^3$;
przy $\chi^{\mathrm{orb}}<0$ geometria
$\widetilde{\mathrm{SL}}_2(\mathbb R)$ lub
$\mathbb H^2\times\mathbb R$, z rozróżnieniem według
zanikania liczby Eulera włóknowania.

\section{Od Ricciego do pełnej krzywizny w wymiarze trzecim}
\label{sec:thurston-ricci}

Wymiar trzy ma jeszcze jedną niezwykłą własność: tensor Riemanna
jest całkowicie określony przez tensor Ricciego. Jest to konsekwencja
zanikania tensora Weyla w wymiarze trzy.

\begin{twierdzenie}[Ricci determinuje krzywiznę w wymiarze trzy]
Niech $(M^3,g)$ będzie trójwymiarową rozmaitością Riemanna.
W każdym punkcie można wybrać ortonormalną bazę
$e_1,e_2,e_3$ diagonalizującą tensor Ricciego, o wartościach
własnych $r_1,r_2,r_3$. Wtedy krzywizny trzech płaszczyzn
własnych są równe
\[
K_{12}=\frac{r_1+r_2-r_3}{2},\quad
K_{13}=\frac{r_1+r_3-r_2}{2},\quad
K_{23}=\frac{r_2+r_3-r_1}{2}.
\]
W szczególności $\operatorname{Ric}=0$ wtedy i tylko wtedy,
gdy pełny tensor krzywizny Riemanna jest zerowy.
\end{twierdzenie}
\begin{proof}
Przypomnijmy, że
$r_1=K_{12}+K_{13}$,
$r_2=K_{12}+K_{23}$ i
$r_3=K_{13}+K_{23}$.
Rozwiązanie tego układu liniowego daje podane wzory.
W wymiarze trzy wszystkie składowe tensora krzywizny można
odzyskać z tensora Ricciego (równoważnie: część Weyla jest
identycznie zerowa). W szczególności zerowanie się wszystkich
wartości własnych tensora Ricciego daje zerowanie się tensora
Riemanna.
\end{proof}

\begin{wniosek}
Jeżeli trójwymiarowa rozmaitość Riemanna spełnia
$\operatorname{Ric}=\lambda g$ dla pewnej stałej $\lambda$,
to ma stałą krzywiznę przekrojową $K=\lambda/2$.
Każda zupełna jednospójna trójrozmaitość spełniająca
$\operatorname{Ric}=0$ jest izometryczna z $\mathbb R^3$.
\end{wniosek}
\begin{proof}
Pierwsza teza wynika bezpośrednio z powyższych wzorów.
W drugim przypadku metryka jest płaska; zastosowanie
klasyfikacji zupełnych jednospójnych przestrzeni o stałej
krzywiźnie daje $\mathbb R^3$.
\end{proof}

\begin{uwaga}[Co zmienia się w wyższych wymiarach?]
Dla $n\ge4$ tensor Weyla może być niezerowy, więc metryka
Ricci-płaska nie musi być płaska. Własność wykorzystana
powyżej jest specyficzna dla wymiaru trzy i wyjaśnia,
dlaczego warunek $\operatorname{Ric}=0$ ma w geometrii
trójwymiarowej znacznie silniejsze konsekwencje niż w
geometrii czterowymiarowej.
\end{uwaga}

\section{Hiperboliczne trójrozmaitości i wielościany fundamentalne}
\label{sec:thurston-hyperbolic}

Konstrukcja zwartych ilorazów $\mathbb H^3$ wymaga innej techniki
niż torus odwzorowania czy kratowe ilorazy grup Liego.
Naturalnym narzędziem są wielościany fundamentalne i izometrie
parujące ich ściany.

\begin{definicja}[Dziedzina fundamentalna]
Niech $\Gamma$ będzie dyskretną grupą izometrii $X$.
Domknięty zbiór $P\subset X$ nazywamy \emph{dziedziną
fundamentalną}, jeżeli $\Gamma P=X$, a wnętrza jego różnych
translat są rozłączne, tzn.
$\operatorname{int}(P)\cap\gamma\operatorname{int}(P)=\varnothing$
dla $\gamma\in\Gamma\setminus\{e\}$.
\end{definicja}

Twierdzenie Poincarégo o wielościanach daje warunki, pod którymi
przyporządkowanie izometrii parom ścian wielościanu definiuje
dyskretną grupę, a sam wielościan jest fundamentalny. W pełnej
wersji należy zweryfikować zgodność parowania ścian,
relacje wokół krawędzi i warunki lokalnego domknięcia oraz
zupełności. Sam fakt, że ściany dają się sparować, nie
wystarcza do otrzymania poprawnej rozmaitości hiperbolicznej.

\begin{przyklad}[Dopełnienie węzła ósemkowego]
Niech $K\subset\mathbb S^3$ będzie węzłem ósemkowym.
Wnętrze dopełnienia $\mathbb S^3\setminus N(K)$, gdzie $N(K)$
jest regularnym otoczeniem węzła, ma zupełną strukturę
hiperboliczną o skończonej objętości. Można ją zbudować przez
sklejenie dwóch regularnych idealnych czworościanów
hiperbolicznych. Jeśli $v_3$ oznacza objętość regularnego
idealnego czworościanu, to objętość tej rozmaitości wynosi
$2v_3$. Nie jest to przykład \emph{zwartej} rozmaitości:
na końcu znajduje się cusp (koniec paraboliczny) z
przekrojem będącym dwuwymiarowym torusem.
\end{przyklad}

\begin{przyklad}[Rozmaitość Giesekinga]
Rozmaitość Giesekinga można otrzymać przez odpowiednie
identyfikacje ścian jednego regularnego idealnego
czworościanu hiperbolicznego. Jest ona nieorientowalna,
niezwarta, zupełna i ma skończoną objętość $v_3$.
Jej orientowalne nakrycie dwukrotne jest izometryczne z
dopełnieniem węzła ósemkowego (z naturalną metryką
hiperboliczną). Ten przykład pokazuje, że topologia końców
i orientowalność są istotne nawet przy stałej krzywiźnie $-1$.
\end{przyklad}

\section{Klasyfikacja lokalnych geometrii a geometryzacja}

Nie należy utożsamiać dwóch twierdzeń: klasyfikacja Thurstona
opisuje osiem możliwych \emph{maksymalnych geometrii modelowych}
na jednospójnych jednorodnych przestrzeniach, natomiast
\emph{geometryzacja} dotyczy globalnej struktury topologicznej
dowolnych trójrozmaitości.

\begin{twierdzenie}[Geometryzacja Thurstona--Perelmana, schemat]
Każda zwarta orientowalna trójrozmaitość po rozkładzie na
składniki pierwsze (cięcia wzdłuż sfer dwuwymiarowych),
a następnie odpowiednim rozkładzie wzdłuż nieściśliwych
torusów, daje części typu Seiferta i części hiperboliczne.
Hiperboliczne wnętrza odpowiednich części mają zupełne
metryki hiperboliczne o skończonej objętości.
Wśród możliwych geometrii składników zamkniętych i
przykładów podstawowych pojawia się osiem geometrii
z twierdzenia~\ref{thm:osiem-geometrii}.
\end{twierdzenie}

Sformułowanie jest celowo schematyczne: poprawna wersja
wymaga uwzględnienia redukcji sferycznej, rozkładu JSJ,
przypadków o pustym rozkładzie i sposobu traktowania brzegu.
Dowód geometryzacji nie jest konsekwencją samej listy ośmiu
geometrii. Historycznie przełomem było zastosowanie przepływu
Ricciego z chirurgią; schemat rozwoju metryki zapisuje
równanie
\[
 \frac{\partial g(t)}{\partial t}=-2\operatorname{Ric}(g(t)).
\]
Jego szczegółowa analiza wymaga oddzielnego rozdziału.

\begin{uwaga}[Różne znaczenia słowa „geometria”]
Nie każda zwarta trójrozmaitość dopuszcza jedną geometrię
jednorodną na całej swojej przestrzeni. Geometryzacja mówi,
że po odpowiednich cięciach otrzymujemy części modelowane
przez geometrie z listy. Jest to zasadnicza różnica między
rozmaitością \emph{geometryczną} a rozmaitością
\emph{geometryzowalną}.
\end{uwaga}

\section{Podsumowanie i powiązania z poprzednimi rozdziałami}

Osiem geometrii Thurstona powstaje z dwóch współdziałających
idei: jednorodności metrycznej i istnienia zwartych ilorazów.
Z punktu widzenia rachunku różniczkowego najważniejsze są
globalne ortonormalne ramy, wzór Koszula i obliczenia krzywizn
Nil oraz Sol. Z punktu widzenia topologii zasadniczą rolę
odgrywają grupy podstawowe, nakrycia, wiązki okręgów, liczba
Eulera i torusy odwzorowania. Geometria
$\widetilde{\mathrm{SL}}_2(\mathbb R)$ łączy oba punkty widzenia
przez krzywiznę połączenia wiązki Seiferta.

Warto przy tym pamiętać, że twierdzenie o geometryzacji
ma charakter globalny i nie wynika z lokalnej listy geometrii.
Jest ono pomostem do bardziej zaawansowanej topologii
trójrozmaitości, teorii rozkładów JSJ i przepływu Ricciego.

\paragraph{Literatura uzupełniająca.}
W.~P.~Thurston, \emph{Three-Dimensional Geometry and Topology},
vol.~1, Princeton University Press, 1997;
P.~Scott, \emph{The Geometries of 3-Manifolds},
Bulletin of the London Mathematical Society 15 (1983), 401--487;
J.~W.~Milnor, \emph{Curvatures of Left Invariant Metrics on Lie Groups},
Advances in Mathematics 21 (1976), 293--329;
S.~A.~Pinillos Prado, \emph{Geometric Structures on 3-Dimensional
Manifolds}, Universidade Federal de Minas Gerais, 2021.

% Wymagane w preambule: amsmath, amssymb, amsthm, mathtools.
% Srodowiska: definicja, twierdzenie, przyklad, uwaga, wniosek, proof.
% Jezeli w skrypcie stosowana jest inna definicja srodowisk,
% nalezy dostosowac nazwy w preambule, nie w tresci rozdzialu.

\chapter{Geometrie Thurstona i struktury geometryczne na trójrozmaitościach}
\label{chap:geometrie-thurstona}

W poprzednich rozdziałach rozważaliśmy rozmaitości Riemanna, ich połączenia,
geodezyjne i krzywiznę oraz podstawowe geometrie przestrzeni o stałej krzywiźnie.
Tutaj przejdziemy od geometrii pojedynczej rozmaitości do pytania o to,
\emph{jakie geometrie jednorodne mogą występować lokalnie na trójrozmaitościach}.
Z jednej strony prowadzi to do klasyfikacji ośmiu geometrii modelowych
Thurstona, z drugiej -- do związków między metryką a topologią rozmaitości
trójwymiarowych. W szczególności zobaczymy, dlaczego do opisu wszystkich
zwartych trójrozmaitości nie wystarczają geometrie sferyczna, euklidesowa
i hiperboliczna.

\section{Struktury \texorpdfstring{$(G,X)$}{(G,X)} i geometrie modelowe}

Ideą programu erlangeńskiego Kleina jest badanie przestrzeni za pomocą
grupy przekształceń zachowujących interesującą nas strukturę. Geometria
Riemanna dostarcza szczególnego przypadku: jeśli $(X,g)$ jest rozmaitością
Riemanna, to grupa $\operatorname{Isom}(X,g)$ składa się z jej izometrii.

\begin{definicja}[Struktura $(G,X)$]
Niech $X$ będzie spójną gładką rozmaitością, a $G$ grupą Liego działającą na
$X$ przez dyfeomorfizmy. \emph{Strukturą $(G,X)$} na rozmaitości $M$
nazywamy atlas map $\{(U_\alpha,\varphi_\alpha)\}$,
gdzie $\varphi_\alpha\colon U_\alpha\to V_\alpha\subset X$ jest
dyfeomorfizmem na zbiór otwarty, a dla każdej spójnej składowej
$C\subset U_\alpha\cap U_\beta$ istnieje $g_{\alpha\beta,C}\in G$ takie, że
$\varphi_\alpha\circ\varphi_\beta^{-1}$ na $\varphi_\beta(C)$
jest obcięciem działania $g_{\alpha\beta,C}$.
\end{definicja}

Warunek dotyczący składowych przecięcia jest istotny: jedno przecięcie
$U_\alpha\cap U_\beta$ może być niespójne, a na różnych jego składowych
mogą występować różne elementy grupy $G$. Zwykły atlas gładki odpowiada
pseudogrupie wszystkich lokalnych dyfeomorfizmów, podczas gdy struktura
$(G,X)$ dopuszcza tylko lokalne przekształcenia pochodzące z jednej grupy.

\begin{przyklad}[Torus płaski]
Niech $X=\mathbb R^3$, $G=\operatorname{Isom}(\mathbb R^3)$ i
$\Gamma=\mathbb Z^3$ działająca przez translacje. Iloraz
$\mathbb T^3=\mathbb R^3/\mathbb Z^3$ otrzymuje atlas $(G,X)$, ponieważ
lokalne odwrotności nakrycia różnią się o translacje należące do $G$.
Wszystkie przestrzenie styczne są euklidesowe, chociaż $\mathbb T^3$
nie jest globalnie izometryczny z $\mathbb R^3$.
\end{przyklad}

\begin{definicja}[Geometria modelowa Thurstona]
\label{def:model-thurston}
Przez \emph{geometrię modelową} rozumiemy parę $(G,X)$, dla której:
\begin{enumerate}
\item $X$ jest spójna i jednospójna;
\item $G$ działa na $X$ tranzytywnie, a stabilizator punktu
      $G_x=\{a\in G:a\cdot x=x\}$ jest zwarty;
\item $G$ jest maksymalną grupą dyfeomorfizmów $X$ spełniającą warunek
      zwartości stabilizatorów;
\item istnieje zwarta rozmaitość z pełną strukturą $(G,X)$.
\end{enumerate}
Dwie geometrie utożsamiamy, gdy ich przestrzenie i grupy działań są
izomorficzne w naturalnym sensie. Przymiotnik \emph{pełna} oznacza, że
struktura pochodzi z ilorazu $X/\Gamma$ dla odpowiednio działającej
grupy dyskretnej $\Gamma\subset G$.
\end{definicja}

Dwa ostatnie warunki nie są kosmetyczne. Gdyby zrezygnować z maksymalności,
tę samą geometrię euklidesową można by sztucznie traktować jako wiele
różnych geometrii, wybierając mniejsze grupy ruchów. Wymóg zwartego
ilorazu eliminuje z kolei liczne jednorodne przestrzenie niemające
zwartych realizacji globalnych.

\begin{twierdzenie}[Metryka niezmiennicza]
\label{thm:metryka-niezmiennicza}
Jeżeli grupa Liego $G$ działa tranzytywnie na $X$, a stabilizator $G_o$
pewnego punktu $o\in X$ jest zwarty, to na $X$ istnieje metryka
Riemanna niezmiennicza względem $G$.
\end{twierdzenie}
\begin{proof}
Wybierzmy dowolny iloczyn skalarny $h$ na $T_oX$. Niech $\mu$ będzie
unormowaną miarą Haara na zwartej grupie $G_o$. Definiujemy
\[
 h_o^{G}(u,v)=\int_{G_o}h(da_o u,da_o v)\,d\mu(a).
\]
Otrzymana forma jest dodatnio określona i niezmiennicza względem
reprezentacji izotropii. Jeśli $x=b\cdot o$, przenosimy ją wzorem
$h_x^G(db_o u,db_o v)=h_o^G(u,v)$. Definicja nie zależy od wyboru $b$:
dwa takie elementy różnią się o element stabilizatora $G_o$, a forma
$h_o^G$ została względem niego uśredniona. Gładkość metryki wynika z
gładkości działania i lokalnego wyboru przekrojów odwzorowania
$G\to G/G_o\simeq X$.
\end{proof}

\begin{uwaga}[Nakrycie uniwersalne i grupa decka]
Jeżeli $X$ jest jednospójna, $\Gamma\subset\operatorname{Isom}(X,g)$
działa swobodnie, właściwie nieciągle i współzwarto, to
$M=\Gamma\backslash X$ jest zwartą rozmaitością Riemanna, a
$X\to M$ jest jej nakryciem uniwersalnym. W szczególności
$\pi_1(M)\simeq\Gamma$. Odwrotnie, każda \emph{pełna} struktura
$(G,X)$ na zwartej rozmaitości jest realizowana w ten sposób.
Nie każda struktura $(G,X)$ jest jednak z definicji pełna.
\end{uwaga}

Przydatne pojęcie pośrednie daje \emph{odwzorowanie rozwijające}
$\operatorname{dev}\colon\widetilde M\to X$. Powstaje ono przez analityczne
kontynuowanie kolejnych lokalnych map modelowych na nakryciu
uniwersalnym i spełnia
$\operatorname{dev}(\gamma p)=\rho(\gamma)\operatorname{dev}(p)$,
gdzie $\rho\colon\pi_1(M)\to G$ jest \emph{holonomią}. Jeżeli struktura
jest pełna, odwzorowanie rozwijające identyfikuje $\widetilde M$ z $X$.

\section{Dlaczego w wymiarze trzecim pojawia się osiem geometrii?}

\begin{twierdzenie}[Klasyfikacja geometrii modelowych Thurstona]
\label{thm:osiem-geometrii}
W wymiarze trzecim, z dokładnością do izomorfizmu geometrii modelowych,
istnieje dokładnie osiem możliwości:
\[
\mathbb S^3,\qquad \mathbb R^3,\qquad \mathbb H^3,\qquad
\mathbb S^2\!\times\!\mathbb R,\qquad
\mathbb H^2\!\times\!\mathbb R,\qquad
\mathrm{Nil},\qquad\widetilde{\mathrm{SL}}_2(\mathbb R),\qquad
\mathrm{Sol}.
\]
Dla trzech pierwszych geometrii spójna izotropia jest trójwymiarowa;
dla następnych czterech -- jednowymiarowa; dla geometrii Sol --
zerowymiarowa.
\end{twierdzenie}

Podamy mechanizm klasyfikacji, oddzielając go od pełnego dowodu, który
wymaga dodatkowej analizy jednorodnych grup Liego i ich sieci.
Na mocy twierdzenia~\ref{thm:metryka-niezmiennicza} możemy wybrać
metrykę niezmienniczą na $X$. Izotropia działa wiernie na $T_oX$
(izometria ustalająca punkt i mająca identyczną różniczkę jest
identycznością), zatem jej składowa spójna jest zwartą spójną
podgrupą $\mathrm{SO}(3)$. Takie podgrupy są, z dokładnością do
sprzężenia, równe $\mathrm{SO}(3)$, $\mathrm{SO}(2)$ albo grupie
trywialnej.

Gdy izotropia ma wymiar trzy, działa tranzytywnie na kierunkach
w każdej przestrzeni stycznej. Krzywizna przekrojowa zależy więc
tylko od punktu, a jednorodność czyni ją stałą na całym $X$.
Otrzymujemy trzy przestrzenie stałej krzywizny.

Gdy izotropia ma wymiar jeden, zachowuje pewien wyróżniony kierunek
pola prostych. Na odpowiednim nakryciu pole to można zorientować i
otrzymać niezmiennicze pole wektorowe, którego linie całkowe tworzą
jednowymiarowe włóknowanie. Geometria poprzeczna jest jedną z trzech
geometrii powierzchni jednorodnych; krzywizna połączenia tego
włóknowania odróżnia geometrie produktowe od Nil i
$\widetilde{\mathrm{SL}}_2(\mathbb R)$. Przy zerowymiarowej izotropii
analizuje się trójwymiarowe grupy Liego wyposażone w metryki
lewostronnie niezmiennicze. Po uwzględnieniu maksymalności grupy
izometrii i istnienia zwartego ilorazu jedyną nową geometrią jest Sol.
To jest \emph{szkic dowodu}, nie pełna klasyfikacja możliwych algebr
Liego i ich sieci.

\begin{center}
\begin{tabular}{|l|c|l|}
\hline
Geometria & $\dim G_o^0$ & Typ geometrii poprzecznej lub krzywizna \\
\hline
$\mathbb S^3$ & $3$ & $K\equiv 1$ \\
$\mathbb R^3$ & $3$ & $K\equiv 0$ \\
$\mathbb H^3$ & $3$ & $K\equiv -1$ \\
$\mathbb S^2\times\mathbb R$ & $1$ & produkt, baza $\mathbb S^2$ \\
$\mathbb H^2\times\mathbb R$ & $1$ & produkt, baza $\mathbb H^2$ \\
Nil & $1$ & baza $\mathbb R^2$, skręcone włókna \\
$\widetilde{\mathrm{SL}}_2(\mathbb R)$ & $1$ & baza $\mathbb H^2$, skręcone włókna \\
Sol & $0$ & włóknowanie nad $\mathbb R$ \\
\hline
\end{tabular}
\end{center}

\section{Trzy geometrie o stałej krzywiźnie}

Przypomnijmy jedynie elementy potrzebne do dalszego porównania. Po
odpowiednim przeskalowaniu metryk geometrie te mają krzywiznę
przekrojową odpowiednio $\kappa=1,0,-1$. W wymiarze trzy ich
tensor Ricciego i krzywizna skalarna spełniają
$\operatorname{Ric}=2\kappa g$ oraz $\operatorname{Scal}=6\kappa$.

\begin{przyklad}[Modele i zwarte ilorazy]
\begin{enumerate}
\item $\mathbb S^3=\{z\in\mathbb C^2:|z_1|^2+|z_2|^2=1\}$,
      z metryką indukowaną z $\mathbb R^4$. Przykładami zwartych
      ilorazów są przestrzenie soczewkowe
      $L(p,q)=\mathbb S^3/\mathbb Z_p$, gdzie generator działa przez
      $(z_1,z_2)\mapsto(e^{2\pi i/p}z_1,e^{2\pi iq/p}z_2)$,
      a $\gcd(p,q)=1$.
\item $\mathbb R^3$, z metryką $g=dx^2+dy^2+dz^2$;
      ilorazem zwartym jest $\mathbb T^3$.
\item $\mathbb H^3=\{(x,y,z)\in\mathbb R^3:z>0\}$,
      z metryką $g=z^{-2}(dx^2+dy^2+dz^2)$. Istnieją zwarte ilorazy
      $\Gamma\backslash\mathbb H^3$, gdzie $\Gamma$ jest bezskrętną
      współzwartą kratą w $\operatorname{Isom}(\mathbb H^3)$.
\end{enumerate}
\end{przyklad}

\section{Dwie geometrie produktowe}

Niech $(\Sigma^2,g_\Sigma)$ będzie powierzchnią o stałej krzywiźnie
Gaussa $\kappa\in\{1,-1\}$. Na $X=\Sigma\times\mathbb R$ bierzemy
metrykę produktową $g=g_\Sigma+dt^2$. Ponieważ połączenie
Levi-Civity produktu zachowuje oba czynniki, dla ortonormalnych
wektorów $e_1,e_2$ stycznych do $\Sigma$ oraz $e_3=\partial_t$
mamy $K(e_1,e_2)=\kappa$ i $K(e_1,e_3)=K(e_2,e_3)=0$.
W konsekwencji
\[
\operatorname{Ric}(e_1,e_1)=\operatorname{Ric}(e_2,e_2)=\kappa,
\qquad \operatorname{Ric}(e_3,e_3)=0,
\qquad \operatorname{Scal}=2\kappa.
\]
Są to geometrie jednorodne, lecz \emph{nie} mają stałej krzywizny
przekrojowej w wymiarze trzy.

\begin{przyklad}[Zwarte realizacje]
Dla $\mathbb S^2\times\mathbb R$ otrzymujemy
$\mathbb S^2\times\mathbb S^1$ przez ilorazowanie drugiego czynnika
translacjami całkowitymi. Dla $\mathbb H^2\times\mathbb R$ wybieramy
zwartą powierzchnię hiperboliczną $\Sigma_g=\Gamma\backslash\mathbb H^2$
rodzaju $g\ge2$; wtedy $\Sigma_g\times\mathbb S^1$ jest zwartym
ilorazem $\mathbb H^2\times\mathbb R$.
\end{przyklad}

\section{Geometria Nil: grupa Heisenberga}
\label{sec:thurston-nil}

Pierwsza istotnie nowa geometria pojawia się wtedy, gdy płaszczyznę
$\mathbb R^2$ zastąpimy rodziną prostych, których połączenie ma
niezerową krzywiznę. Rozważmy $\mathbb R^3$ z działaniem grupowym
\[
(x,y,z)\cdot(x',y',z')=
(x+x',\,y+y',\,z+z'+xy').
\]
Jest to trójwymiarowa \emph{grupa Heisenberga} $H_3$. Można ją
utożsamić z grupą macierzy
$\left(\begin{smallmatrix}1&x&z\\0&1&y\\0&0&1\end{smallmatrix}\right)$.

Jednoformy $\theta^1=dx$, $\theta^2=dy$ i
$\theta^3=dz-x\,dy$ są lewostronnie niezmiennicze.
Ich formami dualnymi są pola
$E_1=\partial_x$, $E_2=\partial_y+x\partial_z$ oraz
$E_3=\partial_z$. Bezpośrednio
$[E_1,E_2]=E_3$, zaś pozostałe nawiasy są zerowe.
Nadajmy $H_3$ metrykę
\begin{equation}\label{eq:nil-metric}
 g_{\mathrm{Nil}}=(\theta^1)^2+(\theta^2)^2+(\theta^3)^2
 =dx^2+dy^2+(dz-x\,dy)^2.
\end{equation}
Zapis $\theta^2$ w powyższej formule oznacza iloczyn tensorowy
$\theta\otimes\theta$, a nie iloczyn zewnętrzny. Pola $E_1,E_2,E_3$
tworzą globalną ortonormalną ramę.

\begin{przyklad}[Pozioma krzywa i holonomia]
Rzut $\pi\colon H_3\to\mathbb R^2$, $\pi(x,y,z)=(x,y)$,
ma włókna równoległe do osi $z$. Rozkład poziomy określa
$\theta^3=0$, czyli dla krzywej $\gamma(s)=(x(s),y(s),z(s))$
warunek poziomości brzmi $z'(s)=x(s)y'(s)$. Po przejściu
wzdłuż zamkniętej pętli $c$ w płaszczyźnie podstawa zostaje
niezmieniona, ale współrzędna $z$ zmienia się o
\[
\Delta z=\oint_c x\,dy=\iint_D dx\wedge dy,
\]
gdzie $D$ jest zorientowanym obszarem ograniczonym przez $c$.
Oznacza to, że przesunięcie wzdłuż włókna mierzy pole zorientowane
ograniczone przez pętlę. Ponieważ $d\theta^3=-dx\wedge dy$,
rozkład poziomy nie jest całkowalny.
\end{przyklad}

\begin{twierdzenie}[Krzywizny geometrii Nil]
Dla metryki~\eqref{eq:nil-metric} krzywizny przekrojowe płaszczyzn
wyznaczonych przez ramę $E_1,E_2,E_3$ wynoszą
$K_{12}=-3/4$ i $K_{13}=K_{23}=1/4$.
Stąd $\operatorname{Ric}=\operatorname{diag}(-1/2,-1/2,1/2)$
w tej ramie i $\operatorname{Scal}=-1/2$.
\end{twierdzenie}
\begin{proof}
Dla ortonormalnych pól lewostronnie niezmienniczych wzór Koszula
przybiera postać
\[
2\langle\nabla_XY,Z\rangle
=\langle[X,Y],Z\rangle-\langle[Y,Z],X\rangle
+\langle[Z,X],Y\rangle.
\]
Z jedynego niezerowego nawiasu $[E_1,E_2]=E_3$ dostajemy
\[
\begin{aligned}
&\nabla_{E_1}E_2=\tfrac12E_3,\qquad
 \nabla_{E_2}E_1=-\tfrac12E_3,\\
&\nabla_{E_1}E_3=\nabla_{E_3}E_1=-\tfrac12E_2,\\
&\nabla_{E_2}E_3=\nabla_{E_3}E_2=\tfrac12E_1,
\qquad \nabla_{E_i}E_i=0.
\end{aligned}
\]
Przy konwencji $R(X,Y)Z=\nabla_X\nabla_YZ-\nabla_Y\nabla_XZ-
\nabla_{[X,Y]}Z$ otrzymujemy na przykład
$R(E_1,E_2)E_2=-\tfrac34E_1$ oraz
$R(E_1,E_3)E_3=\tfrac14E_1$. Trzeci rachunek jest analogiczny.
Tensor Ricciego jest sumą krzywizn przekrojowych po ortonormalnych
kierunkach; daje to podane wartości.
\end{proof}

\begin{przyklad}[Zwarta nilrozmaitość]
Podgrupa $\Gamma=H_3(\mathbb Z)$ macierzy Heisenberga o całkowitych
współrzędnych jest dyskretna i współzwarta. Iloraz
$\Gamma\backslash H_3(\mathbb R)$ jest zwartą trójrozmaitością
z geometrią Nil. Stanowi nietrywialną wiązkę okręgów nad
$\mathbb T^2$: przechodząc wokół podstawowego równoległoboku w bazie,
włókno doznaje przesunięcia wynikającego z niezerowej krzywizny
połączenia. Nie jest to torus płaski $\mathbb T^3$.
\end{przyklad}

\section{Geometria \texorpdfstring{$\widetilde{\mathrm{SL}}_2(\mathbb R)$}{SL2-tylda}}
\label{sec:thurston-sl2}

Geometria ta jest naturalnie związana z jednostkową wiązką styczną
płaszczyzny hiperbolicznej. Grupa izometrii zachowujących orientację
$\mathbb H^2$ jest izomorficzna z
$\operatorname{PSL}_2(\mathbb R)$; działa ona swobodnie i tranzytywnie
na wektorach stycznych długości jeden. Stąd
$T^1\mathbb H^2\simeq\operatorname{PSL}_2(\mathbb R)$ jako rozmaitości.
Nakryciem uniwersalnym tej przestrzeni jest
$\widetilde{\operatorname{SL}}_2(\mathbb R)$.

W modelu górnej półpłaszczyzny $\mathbb H^2=\{(x,y):y>0\}$
można, po identyfikacji nakrycia uniwersalnego z
$\mathbb H^2\times\mathbb R$, wybrać metrykę
\begin{equation}\label{eq:sl2-metric}
 g_{\widetilde{\mathrm{SL}}_2}
 =\frac{dx^2+dy^2}{y^2}
 +\left(dt+\frac{dx}{y}\right)^2.
\end{equation}
Rzut na $(x,y)$ jest submersją Riemannowską o jednowymiarowych,
geodezyjnych włóknach. Forma połączenia
$\alpha=dt+dx/y$ ma krzywiznę
$d\alpha=dx\wedge dy/y^2$, a więc jest ona (z uwzględnieniem
orientacji) równa formie pola powierzchni hiperbolicznej.
Nie jest to metryka produktowa.

\begin{uwaga}[Krzywizny]
Dla normalizacji~\eqref{eq:sl2-metric} wzory O'Neilla dla
submersji Riemannowskiej dają krzywiznę poziomej płaszczyzny
$-1-3/4=-7/4$ oraz krzywizny płaszczyzn zawierających kierunek
włókna równe $1/4$. W szczególności w poziomej ortonormalnej
bazie $E_1,E_2$ i pionowym $E_3$ tensor Ricciego ma wartości
$(-3/2,-3/2,1/2)$. Konkretne wartości zależą od normalizacji
współczynnika przy $\alpha^2$; znak niezerowej krzywizny
połączenia jest ważniejszy od tych liczb.
\end{uwaga}

\begin{przyklad}[Wiązka jednostkowych wektorów nad powierzchnią]
Niech $\Sigma_g=\Gamma\backslash\mathbb H^2$ będzie zamkniętą,
zorientowaną powierzchnią hiperboliczną rodzaju $g\ge2$.
Wówczas $T^1\Sigma_g$ jest zwartą trójrozmaitością, która otrzymuje
geometrię $\widetilde{\operatorname{SL}}_2(\mathbb R)$.
Wiązka $\mathbb S^1\to T^1\Sigma_g\to\Sigma_g$ ma liczbę
Eulera $\chi(\Sigma_g)=2-2g\neq0$. Wynika to z faktu, że
$T^1\Sigma_g$ jest wiązką okręgów stowarzyszoną ze zorientowaną
wiązką styczną $T\Sigma_g$, której klasa Eulera paruje z klasą
fundamentalną do charakterystyki Eulera powierzchni.
\end{przyklad}

Powyższy przykład należy porównać z
$\Sigma_g\times\mathbb S^1$ o geometrii
$\mathbb H^2\times\mathbb R$. Obie rozmaitości mają wiązkę okręgów
nad tą samą powierzchnią, ale w pierwszej liczba Eulera wiązki
jest niezerowa, zaś w drugim przypadku zerowa. Topologia
wiązki wykrywa więc różnicę między dwiema geometriami.

\section{Geometria Sol i odwzorowania sklejenia torusa}
\label{sec:thurston-sol}

Grupę Sol można zrealizować jako $\mathbb R^2\rtimes\mathbb R$,
gdzie $t\in\mathbb R$ działa na $\mathbb R^2$ macierzą
$\operatorname{diag}(e^t,e^{-t})$. Prawo grupowe jest dane wzorem
\[
(x,y,t)\cdot(x',y',t')
=(x+e^t x',\,y+e^{-t}y',\,t+t').
\]
W odróżnieniu od Nil, ta grupa jest rozwiązalna, ale nie nilpotentna.
Jednoformy $\theta^1=e^{-t}dx$, $\theta^2=e^t dy$ i
$\theta^3=dt$ są lewostronnie niezmiennicze. Wybieramy metrykę
\begin{equation}\label{eq:sol-metric}
 g_{\mathrm{Sol}}=e^{-2t}dx^2+e^{2t}dy^2+dt^2.
\end{equation}
Jej globalna ortonormalna rama to
$E_1=e^t\partial_x$, $E_2=e^{-t}\partial_y$,
$E_3=\partial_t$; spełnia ona
$[E_3,E_1]=E_1$, $[E_3,E_2]=-E_2$ oraz $[E_1,E_2]=0$.

\begin{twierdzenie}[Krzywizny geometrii Sol]
Dla metryki~\eqref{eq:sol-metric} zachodzi
$K_{12}=1$ oraz $K_{13}=K_{23}=-1$.
W szczególności
$\operatorname{Ric}=\operatorname{diag}(0,0,-2)$
w ramie $E_1,E_2,E_3$ i $\operatorname{Scal}=-2$.
\end{twierdzenie}
\begin{proof}
Ze wzoru Koszula otrzymujemy
$\nabla_{E_1}E_1=E_3$,
$\nabla_{E_1}E_3=-E_1$,
$\nabla_{E_2}E_2=-E_3$,
$\nabla_{E_2}E_3=E_2$;
pozostałe pochodne kowariantne pól ramy są zerowe.
Stąd $R(E_1,E_2)E_2=E_1$ oraz
$R(E_1,E_3)E_3=-E_1$ i $R(E_2,E_3)E_3=-E_2$.
Pozostałe wzory wynikają z definicji krzywizny Ricciego.
\end{proof}

\begin{przyklad}[Zwarta geometria Sol jako torus odwzorowania]
\label{ex:sol-mapping-torus}
Niech
$A=\left(\begin{smallmatrix}2&1\\1&1\end{smallmatrix}\right)
\in\mathrm{SL}_2(\mathbb Z)$. Jej wartości własne są rzeczywiste,
dodatnie i równe
$\lambda=(3+\sqrt5)/2>1$ oraz $\lambda^{-1}$.
Rozważmy \emph{torus odwzorowania}
\[
M_A=\frac{\mathbb T^2\times[0,1]}
{(u,1)\sim(Au,0)}.
\]
Jest to zwarta trójrozmaitość będąca wiązką torusów nad
$\mathbb S^1$. Aby zobaczyć geometrię Sol, wybierzmy liniowe
współrzędne na $\mathbb R^2$ diagonalizujące $A$.
W tych współrzędnych działanie $A$ ma postać
$(x,y)\mapsto(\lambda x,\lambda^{-1}y)$, natomiast
$\mathbb Z^2$ przechodzi w pewną kratę $\Lambda\subset\mathbb R^2$,
niezmienniczą względem tego przekształcenia.
Grupa generowana przez translacje o wektory $\Lambda$ oraz
\[
 (x,y,t)\longmapsto
 (\lambda x,\lambda^{-1}y,t+\log\lambda)
\]
działa na $(\mathrm{Sol},g_{\mathrm{Sol}})$ izometrycznie,
swobodnie i współzwarto. Otrzymany iloraz jest torusem
odwzorowania z monodromią $A$ (ewentualnie $A^{-1}$,
zależnie od konwencji sklejania końców przedziału).
Zatem $M_A$ ma geometrię Sol.
\end{przyklad}

Widać tu bezpośrednio związek algebry z geometrią: hiperboliczne
odwzorowanie torusa ma dwa kierunki własne, z których jeden ulega
rozciąganiu, a drugi ściskaniu. Metryka Sol dokładnie kompensuje te
zmiany przy przesunięciu w kierunku trzeciej współrzędnej.

\section{Wiązki Seiferta a geometrie modelowe}

\begin{definicja}[Wiązka Seiferta]
W uproszczeniu \emph{rozmaitością włóknistą Seiferta} nazywamy
trójrozmaitość rozłożoną na okręgi tak, że każde włókno ma
otoczenie modelowane na odpowiednim włóknistym torusie pełnym.
Włókna wyjątkowe dopuszczają skończone lokalne stabilizatory;
przestrzeń ilorazowa jest wówczas dwuwymiarowym orbifoldem.
Gdy włókien wyjątkowych brak, mamy zwykłą wiązkę okręgów nad
powierzchnią.
\end{definicja}

Przykłady $\mathbb S^2\times\mathbb S^1$,
$\Sigma_g\times\mathbb S^1$, zwarta nilrozmaitość oraz
$T^1\Sigma_g$ pokazują wspólny wzorzec: większość geometrii
Thurstona z jednowymiarową izotropią występuje naturalnie na
rozmaitościach z włóknowaniem przez okręgi.
Odpowiednia geometria zależy od krzywizny orbifoldu bazy
oraz liczby Eulera włóknowania. Dla zorientowanej wiązki
okręgów nad zamkniętą zorientowaną powierzchnią sytuację
ilustruje poniższe zestawienie; w ogólności trzeba uwzględnić
również włókna wyjątkowe.

\begin{center}
\begin{tabular}{|c|c|l|}
\hline
Krzywizna bazy & Liczba Eulera wiązki & Geometria \\
\hline
$+1$ & $0$ & $\mathbb S^2\times\mathbb R$ \\
$0$ & $0$ & $\mathbb R^3$ \\
$0$ & $\neq0$ & Nil \\
$-1$ & $0$ & $\mathbb H^2\times\mathbb R$ \\
$-1$ & $\neq0$ & $\widetilde{\mathrm{SL}}_2(\mathbb R)$ \\
\hline
\end{tabular}
\end{center}

Tabela nie obejmuje sferycznych przestrzeni z niezerową klasą
Eulera, które również bywają włókniste Seiferta. Pełne kryteria
dla orbifoldów wykorzystują ich orbifoldową charakterystykę
Eulera oraz odpowiednio znormalizowaną liczbę Eulera
włóknowania. W wersji ogólnej dla zamkniętych, orientowalnych
rozmaitości Seiferta przy $\chi^{\mathrm{orb}}>0$ występuje
geometria $\mathbb S^3$ lub $\mathbb S^2\times\mathbb R$;
przy $\chi^{\mathrm{orb}}=0$ geometria Nil lub $\mathbb R^3$;
przy $\chi^{\mathrm{orb}}<0$ geometria
$\widetilde{\mathrm{SL}}_2(\mathbb R)$ lub
$\mathbb H^2\times\mathbb R$, z rozróżnieniem według
zanikania liczby Eulera włóknowania.

\section{Od Ricciego do pełnej krzywizny w wymiarze trzecim}
\label{sec:thurston-ricci}

Wymiar trzy ma jeszcze jedną niezwykłą własność: tensor Riemanna
jest całkowicie określony przez tensor Ricciego. Jest to konsekwencja
zanikania tensora Weyla w wymiarze trzy.

\begin{twierdzenie}[Ricci determinuje krzywiznę w wymiarze trzy]
Niech $(M^3,g)$ będzie trójwymiarową rozmaitością Riemanna.
W każdym punkcie można wybrać ortonormalną bazę
$e_1,e_2,e_3$ diagonalizującą tensor Ricciego, o wartościach
własnych $r_1,r_2,r_3$. Wtedy krzywizny trzech płaszczyzn
własnych są równe
\[
K_{12}=\frac{r_1+r_2-r_3}{2},\quad
K_{13}=\frac{r_1+r_3-r_2}{2},\quad
K_{23}=\frac{r_2+r_3-r_1}{2}.
\]
W szczególności $\operatorname{Ric}=0$ wtedy i tylko wtedy,
gdy pełny tensor krzywizny Riemanna jest zerowy.
\end{twierdzenie}
\begin{proof}
Przypomnijmy, że
$r_1=K_{12}+K_{13}$,
$r_2=K_{12}+K_{23}$ i
$r_3=K_{13}+K_{23}$.
Rozwiązanie tego układu liniowego daje podane wzory.
W wymiarze trzy wszystkie składowe tensora krzywizny można
odzyskać z tensora Ricciego (równoważnie: część Weyla jest
identycznie zerowa). W szczególności zerowanie się wszystkich
wartości własnych tensora Ricciego daje zerowanie się tensora
Riemanna.
\end{proof}

\begin{wniosek}
Jeżeli trójwymiarowa rozmaitość Riemanna spełnia
$\operatorname{Ric}=\lambda g$ dla pewnej stałej $\lambda$,
to ma stałą krzywiznę przekrojową $K=\lambda/2$.
Każda zupełna jednospójna trójrozmaitość spełniająca
$\operatorname{Ric}=0$ jest izometryczna z $\mathbb R^3$.
\end{wniosek}
\begin{proof}
Pierwsza teza wynika bezpośrednio z powyższych wzorów.
W drugim przypadku metryka jest płaska; zastosowanie
klasyfikacji zupełnych jednospójnych przestrzeni o stałej
krzywiźnie daje $\mathbb R^3$.
\end{proof}

\begin{uwaga}[Co zmienia się w wyższych wymiarach?]
Dla $n\ge4$ tensor Weyla może być niezerowy, więc metryka
Ricci-płaska nie musi być płaska. Własność wykorzystana
powyżej jest specyficzna dla wymiaru trzy i wyjaśnia,
dlaczego warunek $\operatorname{Ric}=0$ ma w geometrii
trójwymiarowej znacznie silniejsze konsekwencje niż w
geometrii czterowymiarowej.
\end{uwaga}

\section{Hiperboliczne trójrozmaitości i wielościany fundamentalne}
\label{sec:thurston-hyperbolic}

Konstrukcja zwartych ilorazów $\mathbb H^3$ wymaga innej techniki
niż torus odwzorowania czy kratowe ilorazy grup Liego.
Naturalnym narzędziem są wielościany fundamentalne i izometrie
parujące ich ściany.

\begin{definicja}[Dziedzina fundamentalna]
Niech $\Gamma$ będzie dyskretną grupą izometrii $X$.
Domknięty zbiór $P\subset X$ nazywamy \emph{dziedziną
fundamentalną}, jeżeli $\Gamma P=X$, a wnętrza jego różnych
translat są rozłączne, tzn.
$\operatorname{int}(P)\cap\gamma\operatorname{int}(P)=\varnothing$
dla $\gamma\in\Gamma\setminus\{e\}$.
\end{definicja}

Twierdzenie Poincarégo o wielościanach daje warunki, pod którymi
przyporządkowanie izometrii parom ścian wielościanu definiuje
dyskretną grupę, a sam wielościan jest fundamentalny. W pełnej
wersji należy zweryfikować zgodność parowania ścian,
relacje wokół krawędzi i warunki lokalnego domknięcia oraz
zupełności. Sam fakt, że ściany dają się sparować, nie
wystarcza do otrzymania poprawnej rozmaitości hiperbolicznej.

\begin{przyklad}[Dopełnienie węzła ósemkowego]
Niech $K\subset\mathbb S^3$ będzie węzłem ósemkowym.
Wnętrze dopełnienia $\mathbb S^3\setminus N(K)$, gdzie $N(K)$
jest regularnym otoczeniem węzła, ma zupełną strukturę
hiperboliczną o skończonej objętości. Można ją zbudować przez
sklejenie dwóch regularnych idealnych czworościanów
hiperbolicznych. Jeśli $v_3$ oznacza objętość regularnego
idealnego czworościanu, to objętość tej rozmaitości wynosi
$2v_3$. Nie jest to przykład \emph{zwartej} rozmaitości:
na końcu znajduje się cusp (koniec paraboliczny) z
przekrojem będącym dwuwymiarowym torusem.
\end{przyklad}

\begin{przyklad}[Rozmaitość Giesekinga]
Rozmaitość Giesekinga można otrzymać przez odpowiednie
identyfikacje ścian jednego regularnego idealnego
czworościanu hiperbolicznego. Jest ona nieorientowalna,
niezwarta, zupełna i ma skończoną objętość $v_3$.
Jej orientowalne nakrycie dwukrotne jest izometryczne z
dopełnieniem węzła ósemkowego (z naturalną metryką
hiperboliczną). Ten przykład pokazuje, że topologia końców
i orientowalność są istotne nawet przy stałej krzywiźnie $-1$.
\end{przyklad}

\section{Klasyfikacja lokalnych geometrii a geometryzacja}

Nie należy utożsamiać dwóch twierdzeń: klasyfikacja Thurstona
opisuje osiem możliwych \emph{maksymalnych geometrii modelowych}
na jednospójnych jednorodnych przestrzeniach, natomiast
\emph{geometryzacja} dotyczy globalnej struktury topologicznej
dowolnych trójrozmaitości.

\begin{twierdzenie}[Geometryzacja Thurstona--Perelmana, schemat]
Każda zwarta orientowalna trójrozmaitość po rozkładzie na
składniki pierwsze (cięcia wzdłuż sfer dwuwymiarowych),
a następnie odpowiednim rozkładzie wzdłuż nieściśliwych
torusów, daje części typu Seiferta i części hiperboliczne.
Hiperboliczne wnętrza odpowiednich części mają zupełne
metryki hiperboliczne o skończonej objętości.
Wśród możliwych geometrii składników zamkniętych i
przykładów podstawowych pojawia się osiem geometrii
z twierdzenia~\ref{thm:osiem-geometrii}.
\end{twierdzenie}

Sformułowanie jest celowo schematyczne: poprawna wersja
wymaga uwzględnienia redukcji sferycznej, rozkładu JSJ,
przypadków o pustym rozkładzie i sposobu traktowania brzegu.
Dowód geometryzacji nie jest konsekwencją samej listy ośmiu
geometrii. Historycznie przełomem było zastosowanie przepływu
Ricciego z chirurgią; schemat rozwoju metryki zapisuje
równanie
\[
 \frac{\partial g(t)}{\partial t}=-2\operatorname{Ric}(g(t)).
\]
Jego szczegółowa analiza wymaga oddzielnego rozdziału.

\begin{uwaga}[Różne znaczenia słowa „geometria”]
Nie każda zwarta trójrozmaitość dopuszcza jedną geometrię
jednorodną na całej swojej przestrzeni. Geometryzacja mówi,
że po odpowiednich cięciach otrzymujemy części modelowane
przez geometrie z listy. Jest to zasadnicza różnica między
rozmaitością \emph{geometryczną} a rozmaitością
\emph{geometryzowalną}.
\end{uwaga}

\section{Podsumowanie i powiązania z poprzednimi rozdziałami}

Osiem geometrii Thurstona powstaje z dwóch współdziałających
idei: jednorodności metrycznej i istnienia zwartych ilorazów.
Z punktu widzenia rachunku różniczkowego najważniejsze są
globalne ortonormalne ramy, wzór Koszula i obliczenia krzywizn
Nil oraz Sol. Z punktu widzenia topologii zasadniczą rolę
odgrywają grupy podstawowe, nakrycia, wiązki okręgów, liczba
Eulera i torusy odwzorowania. Geometria
$\widetilde{\mathrm{SL}}_2(\mathbb R)$ łączy oba punkty widzenia
przez krzywiznę połączenia wiązki Seiferta.

Warto przy tym pamiętać, że twierdzenie o geometryzacji
ma charakter globalny i nie wynika z lokalnej listy geometrii.
Jest ono pomostem do bardziej zaawansowanej topologii
trójrozmaitości, teorii rozkładów JSJ i przepływu Ricciego.

\paragraph{Literatura uzupełniająca.}
W.~P.~Thurston, \emph{Three-Dimensional Geometry and Topology},
vol.~1, Princeton University Press, 1997;
P.~Scott, \emph{The Geometries of 3-Manifolds},
Bulletin of the London Mathematical Society 15 (1983), 401--487;
J.~W.~Milnor, \emph{Curvatures of Left Invariant Metrics on Lie Groups},
Advances in Mathematics 21 (1976), 293--329;
S.~A.~Pinillos Prado, \emph{Geometric Structures on 3-Dimensional
Manifolds}, Universidade Federal de Minas Gerais, 2021.

% Wymagane w preambule: amsmath, amssymb, amsthm, mathtools.
% Srodowiska: definicja, twierdzenie, przyklad, uwaga, wniosek, proof.
% Jezeli w skrypcie stosowana jest inna definicja srodowisk,
% nalezy dostosowac nazwy w preambule, nie w tresci rozdzialu.

\chapter{Geometrie Thurstona i struktury geometryczne na trójrozmaitościach}
\label{chap:geometrie-thurstona}

W poprzednich rozdziałach rozważaliśmy rozmaitości Riemanna, ich połączenia,
geodezyjne i krzywiznę oraz podstawowe geometrie przestrzeni o stałej krzywiźnie.
Tutaj przejdziemy od geometrii pojedynczej rozmaitości do pytania o to,
\emph{jakie geometrie jednorodne mogą występować lokalnie na trójrozmaitościach}.
Z jednej strony prowadzi to do klasyfikacji ośmiu geometrii modelowych
Thurstona, z drugiej -- do związków między metryką a topologią rozmaitości
trójwymiarowych. W szczególności zobaczymy, dlaczego do opisu wszystkich
zwartych trójrozmaitości nie wystarczają geometrie sferyczna, euklidesowa
i hiperboliczna.

\section{Struktury \texorpdfstring{$(G,X)$}{(G,X)} i geometrie modelowe}

Ideą programu erlangeńskiego Kleina jest badanie przestrzeni za pomocą
grupy przekształceń zachowujących interesującą nas strukturę. Geometria
Riemanna dostarcza szczególnego przypadku: jeśli $(X,g)$ jest rozmaitością
Riemanna, to grupa $\operatorname{Isom}(X,g)$ składa się z jej izometrii.

\begin{definicja}[Struktura $(G,X)$]
Niech $X$ będzie spójną gładką rozmaitością, a $G$ grupą Liego działającą na
$X$ przez dyfeomorfizmy. \emph{Strukturą $(G,X)$} na rozmaitości $M$
nazywamy atlas map $\{(U_\alpha,\varphi_\alpha)\}$,
gdzie $\varphi_\alpha\colon U_\alpha\to V_\alpha\subset X$ jest
dyfeomorfizmem na zbiór otwarty, a dla każdej spójnej składowej
$C\subset U_\alpha\cap U_\beta$ istnieje $g_{\alpha\beta,C}\in G$ takie, że
$\varphi_\alpha\circ\varphi_\beta^{-1}$ na $\varphi_\beta(C)$
jest obcięciem działania $g_{\alpha\beta,C}$.
\end{definicja}

Warunek dotyczący składowych przecięcia jest istotny: jedno przecięcie
$U_\alpha\cap U_\beta$ może być niespójne, a na różnych jego składowych
mogą występować różne elementy grupy $G$. Zwykły atlas gładki odpowiada
pseudogrupie wszystkich lokalnych dyfeomorfizmów, podczas gdy struktura
$(G,X)$ dopuszcza tylko lokalne przekształcenia pochodzące z jednej grupy.

\begin{przyklad}[Torus płaski]
Niech $X=\mathbb R^3$, $G=\operatorname{Isom}(\mathbb R^3)$ i
$\Gamma=\mathbb Z^3$ działająca przez translacje. Iloraz
$\mathbb T^3=\mathbb R^3/\mathbb Z^3$ otrzymuje atlas $(G,X)$, ponieważ
lokalne odwrotności nakrycia różnią się o translacje należące do $G$.
Wszystkie przestrzenie styczne są euklidesowe, chociaż $\mathbb T^3$
nie jest globalnie izometryczny z $\mathbb R^3$.
\end{przyklad}

\begin{definicja}[Geometria modelowa Thurstona]
\label{def:model-thurston}
Przez \emph{geometrię modelową} rozumiemy parę $(G,X)$, dla której:
\begin{enumerate}
\item $X$ jest spójna i jednospójna;
\item $G$ działa na $X$ tranzytywnie, a stabilizator punktu
      $G_x=\{a\in G:a\cdot x=x\}$ jest zwarty;
\item $G$ jest maksymalną grupą dyfeomorfizmów $X$ spełniającą warunek
      zwartości stabilizatorów;
\item istnieje zwarta rozmaitość z pełną strukturą $(G,X)$.
\end{enumerate}
Dwie geometrie utożsamiamy, gdy ich przestrzenie i grupy działań są
izomorficzne w naturalnym sensie. Przymiotnik \emph{pełna} oznacza, że
struktura pochodzi z ilorazu $X/\Gamma$ dla odpowiednio działającej
grupy dyskretnej $\Gamma\subset G$.
\end{definicja}

Dwa ostatnie warunki nie są kosmetyczne. Gdyby zrezygnować z maksymalności,
tę samą geometrię euklidesową można by sztucznie traktować jako wiele
różnych geometrii, wybierając mniejsze grupy ruchów. Wymóg zwartego
ilorazu eliminuje z kolei liczne jednorodne przestrzenie niemające
zwartych realizacji globalnych.

\begin{twierdzenie}[Metryka niezmiennicza]
\label{thm:metryka-niezmiennicza}
Jeżeli grupa Liego $G$ działa tranzytywnie na $X$, a stabilizator $G_o$
pewnego punktu $o\in X$ jest zwarty, to na $X$ istnieje metryka
Riemanna niezmiennicza względem $G$.
\end{twierdzenie}
\begin{proof}
Wybierzmy dowolny iloczyn skalarny $h$ na $T_oX$. Niech $\mu$ będzie
unormowaną miarą Haara na zwartej grupie $G_o$. Definiujemy
\[
 h_o^{G}(u,v)=\int_{G_o}h(da_o u,da_o v)\,d\mu(a).
\]
Otrzymana forma jest dodatnio określona i niezmiennicza względem
reprezentacji izotropii. Jeśli $x=b\cdot o$, przenosimy ją wzorem
$h_x^G(db_o u,db_o v)=h_o^G(u,v)$. Definicja nie zależy od wyboru $b$:
dwa takie elementy różnią się o element stabilizatora $G_o$, a forma
$h_o^G$ została względem niego uśredniona. Gładkość metryki wynika z
gładkości działania i lokalnego wyboru przekrojów odwzorowania
$G\to G/G_o\simeq X$.
\end{proof}

\begin{uwaga}[Nakrycie uniwersalne i grupa decka]
Jeżeli $X$ jest jednospójna, $\Gamma\subset\operatorname{Isom}(X,g)$
działa swobodnie, właściwie nieciągle i współzwarto, to
$M=\Gamma\backslash X$ jest zwartą rozmaitością Riemanna, a
$X\to M$ jest jej nakryciem uniwersalnym. W szczególności
$\pi_1(M)\simeq\Gamma$. Odwrotnie, każda \emph{pełna} struktura
$(G,X)$ na zwartej rozmaitości jest realizowana w ten sposób.
Nie każda struktura $(G,X)$ jest jednak z definicji pełna.
\end{uwaga}

Przydatne pojęcie pośrednie daje \emph{odwzorowanie rozwijające}
$\operatorname{dev}\colon\widetilde M\to X$. Powstaje ono przez analityczne
kontynuowanie kolejnych lokalnych map modelowych na nakryciu
uniwersalnym i spełnia
$\operatorname{dev}(\gamma p)=\rho(\gamma)\operatorname{dev}(p)$,
gdzie $\rho\colon\pi_1(M)\to G$ jest \emph{holonomią}. Jeżeli struktura
jest pełna, odwzorowanie rozwijające identyfikuje $\widetilde M$ z $X$.

\section{Dlaczego w wymiarze trzecim pojawia się osiem geometrii?}

\begin{twierdzenie}[Klasyfikacja geometrii modelowych Thurstona]
\label{thm:osiem-geometrii}
W wymiarze trzecim, z dokładnością do izomorfizmu geometrii modelowych,
istnieje dokładnie osiem możliwości:
\[
\mathbb S^3,\qquad \mathbb R^3,\qquad \mathbb H^3,\qquad
\mathbb S^2\!\times\!\mathbb R,\qquad
\mathbb H^2\!\times\!\mathbb R,\qquad
\mathrm{Nil},\qquad\widetilde{\mathrm{SL}}_2(\mathbb R),\qquad
\mathrm{Sol}.
\]
Dla trzech pierwszych geometrii spójna izotropia jest trójwymiarowa;
dla następnych czterech -- jednowymiarowa; dla geometrii Sol --
zerowymiarowa.
\end{twierdzenie}

Podamy mechanizm klasyfikacji, oddzielając go od pełnego dowodu, który
wymaga dodatkowej analizy jednorodnych grup Liego i ich sieci.
Na mocy twierdzenia~\ref{thm:metryka-niezmiennicza} możemy wybrać
metrykę niezmienniczą na $X$. Izotropia działa wiernie na $T_oX$
(izometria ustalająca punkt i mająca identyczną różniczkę jest
identycznością), zatem jej składowa spójna jest zwartą spójną
podgrupą $\mathrm{SO}(3)$. Takie podgrupy są, z dokładnością do
sprzężenia, równe $\mathrm{SO}(3)$, $\mathrm{SO}(2)$ albo grupie
trywialnej.

Gdy izotropia ma wymiar trzy, działa tranzytywnie na kierunkach
w każdej przestrzeni stycznej. Krzywizna przekrojowa zależy więc
tylko od punktu, a jednorodność czyni ją stałą na całym $X$.
Otrzymujemy trzy przestrzenie stałej krzywizny.

Gdy izotropia ma wymiar jeden, zachowuje pewien wyróżniony kierunek
pola prostych. Na odpowiednim nakryciu pole to można zorientować i
otrzymać niezmiennicze pole wektorowe, którego linie całkowe tworzą
jednowymiarowe włóknowanie. Geometria poprzeczna jest jedną z trzech
geometrii powierzchni jednorodnych; krzywizna połączenia tego
włóknowania odróżnia geometrie produktowe od Nil i
$\widetilde{\mathrm{SL}}_2(\mathbb R)$. Przy zerowymiarowej izotropii
analizuje się trójwymiarowe grupy Liego wyposażone w metryki
lewostronnie niezmiennicze. Po uwzględnieniu maksymalności grupy
izometrii i istnienia zwartego ilorazu jedyną nową geometrią jest Sol.
To jest \emph{szkic dowodu}, nie pełna klasyfikacja możliwych algebr
Liego i ich sieci.

\begin{center}
\begin{tabular}{|l|c|l|}
\hline
Geometria & $\dim G_o^0$ & Typ geometrii poprzecznej lub krzywizna \\
\hline
$\mathbb S^3$ & $3$ & $K\equiv 1$ \\
$\mathbb R^3$ & $3$ & $K\equiv 0$ \\
$\mathbb H^3$ & $3$ & $K\equiv -1$ \\
$\mathbb S^2\times\mathbb R$ & $1$ & produkt, baza $\mathbb S^2$ \\
$\mathbb H^2\times\mathbb R$ & $1$ & produkt, baza $\mathbb H^2$ \\
Nil & $1$ & baza $\mathbb R^2$, skręcone włókna \\
$\widetilde{\mathrm{SL}}_2(\mathbb R)$ & $1$ & baza $\mathbb H^2$, skręcone włókna \\
Sol & $0$ & włóknowanie nad $\mathbb R$ \\
\hline
\end{tabular}
\end{center}

\section{Trzy geometrie o stałej krzywiźnie}

Przypomnijmy jedynie elementy potrzebne do dalszego porównania. Po
odpowiednim przeskalowaniu metryk geometrie te mają krzywiznę
przekrojową odpowiednio $\kappa=1,0,-1$. W wymiarze trzy ich
tensor Ricciego i krzywizna skalarna spełniają
$\operatorname{Ric}=2\kappa g$ oraz $\operatorname{Scal}=6\kappa$.

\begin{przyklad}[Modele i zwarte ilorazy]
\begin{enumerate}
\item $\mathbb S^3=\{z\in\mathbb C^2:|z_1|^2+|z_2|^2=1\}$,
      z metryką indukowaną z $\mathbb R^4$. Przykładami zwartych
      ilorazów są przestrzenie soczewkowe
      $L(p,q)=\mathbb S^3/\mathbb Z_p$, gdzie generator działa przez
      $(z_1,z_2)\mapsto(e^{2\pi i/p}z_1,e^{2\pi iq/p}z_2)$,
      a $\gcd(p,q)=1$.
\item $\mathbb R^3$, z metryką $g=dx^2+dy^2+dz^2$;
      ilorazem zwartym jest $\mathbb T^3$.
\item $\mathbb H^3=\{(x,y,z)\in\mathbb R^3:z>0\}$,
      z metryką $g=z^{-2}(dx^2+dy^2+dz^2)$. Istnieją zwarte ilorazy
      $\Gamma\backslash\mathbb H^3$, gdzie $\Gamma$ jest bezskrętną
      współzwartą kratą w $\operatorname{Isom}(\mathbb H^3)$.
\end{enumerate}
\end{przyklad}

\section{Dwie geometrie produktowe}

Niech $(\Sigma^2,g_\Sigma)$ będzie powierzchnią o stałej krzywiźnie
Gaussa $\kappa\in\{1,-1\}$. Na $X=\Sigma\times\mathbb R$ bierzemy
metrykę produktową $g=g_\Sigma+dt^2$. Ponieważ połączenie
Levi-Civity produktu zachowuje oba czynniki, dla ortonormalnych
wektorów $e_1,e_2$ stycznych do $\Sigma$ oraz $e_3=\partial_t$
mamy $K(e_1,e_2)=\kappa$ i $K(e_1,e_3)=K(e_2,e_3)=0$.
W konsekwencji
\[
\operatorname{Ric}(e_1,e_1)=\operatorname{Ric}(e_2,e_2)=\kappa,
\qquad \operatorname{Ric}(e_3,e_3)=0,
\qquad \operatorname{Scal}=2\kappa.
\]
Są to geometrie jednorodne, lecz \emph{nie} mają stałej krzywizny
przekrojowej w wymiarze trzy.

\begin{przyklad}[Zwarte realizacje]
Dla $\mathbb S^2\times\mathbb R$ otrzymujemy
$\mathbb S^2\times\mathbb S^1$ przez ilorazowanie drugiego czynnika
translacjami całkowitymi. Dla $\mathbb H^2\times\mathbb R$ wybieramy
zwartą powierzchnię hiperboliczną $\Sigma_g=\Gamma\backslash\mathbb H^2$
rodzaju $g\ge2$; wtedy $\Sigma_g\times\mathbb S^1$ jest zwartym
ilorazem $\mathbb H^2\times\mathbb R$.
\end{przyklad}

\section{Geometria Nil: grupa Heisenberga}
\label{sec:thurston-nil}

Pierwsza istotnie nowa geometria pojawia się wtedy, gdy płaszczyznę
$\mathbb R^2$ zastąpimy rodziną prostych, których połączenie ma
niezerową krzywiznę. Rozważmy $\mathbb R^3$ z działaniem grupowym
\[
(x,y,z)\cdot(x',y',z')=
(x+x',\,y+y',\,z+z'+xy').
\]
Jest to trójwymiarowa \emph{grupa Heisenberga} $H_3$. Można ją
utożsamić z grupą macierzy
$\left(\begin{smallmatrix}1&x&z\\0&1&y\\0&0&1\end{smallmatrix}\right)$.

Jednoformy $\theta^1=dx$, $\theta^2=dy$ i
$\theta^3=dz-x\,dy$ są lewostronnie niezmiennicze.
Ich formami dualnymi są pola
$E_1=\partial_x$, $E_2=\partial_y+x\partial_z$ oraz
$E_3=\partial_z$. Bezpośrednio
$[E_1,E_2]=E_3$, zaś pozostałe nawiasy są zerowe.
Nadajmy $H_3$ metrykę
\begin{equation}\label{eq:nil-metric}
 g_{\mathrm{Nil}}=(\theta^1)^2+(\theta^2)^2+(\theta^3)^2
 =dx^2+dy^2+(dz-x\,dy)^2.
\end{equation}
Zapis $\theta^2$ w powyższej formule oznacza iloczyn tensorowy
$\theta\otimes\theta$, a nie iloczyn zewnętrzny. Pola $E_1,E_2,E_3$
tworzą globalną ortonormalną ramę.

\begin{przyklad}[Pozioma krzywa i holonomia]
Rzut $\pi\colon H_3\to\mathbb R^2$, $\pi(x,y,z)=(x,y)$,
ma włókna równoległe do osi $z$. Rozkład poziomy określa
$\theta^3=0$, czyli dla krzywej $\gamma(s)=(x(s),y(s),z(s))$
warunek poziomości brzmi $z'(s)=x(s)y'(s)$. Po przejściu
wzdłuż zamkniętej pętli $c$ w płaszczyźnie podstawa zostaje
niezmieniona, ale współrzędna $z$ zmienia się o
\[
\Delta z=\oint_c x\,dy=\iint_D dx\wedge dy,
\]
gdzie $D$ jest zorientowanym obszarem ograniczonym przez $c$.
Oznacza to, że przesunięcie wzdłuż włókna mierzy pole zorientowane
ograniczone przez pętlę. Ponieważ $d\theta^3=-dx\wedge dy$,
rozkład poziomy nie jest całkowalny.
\end{przyklad}

\begin{twierdzenie}[Krzywizny geometrii Nil]
Dla metryki~\eqref{eq:nil-metric} krzywizny przekrojowe płaszczyzn
wyznaczonych przez ramę $E_1,E_2,E_3$ wynoszą
$K_{12}=-3/4$ i $K_{13}=K_{23}=1/4$.
Stąd $\operatorname{Ric}=\operatorname{diag}(-1/2,-1/2,1/2)$
w tej ramie i $\operatorname{Scal}=-1/2$.
\end{twierdzenie}
\begin{proof}
Dla ortonormalnych pól lewostronnie niezmienniczych wzór Koszula
przybiera postać
\[
2\langle\nabla_XY,Z\rangle
=\langle[X,Y],Z\rangle-\langle[Y,Z],X\rangle
+\langle[Z,X],Y\rangle.
\]
Z jedynego niezerowego nawiasu $[E_1,E_2]=E_3$ dostajemy
\[
\begin{aligned}
&\nabla_{E_1}E_2=\tfrac12E_3,\qquad
 \nabla_{E_2}E_1=-\tfrac12E_3,\\
&\nabla_{E_1}E_3=\nabla_{E_3}E_1=-\tfrac12E_2,\\
&\nabla_{E_2}E_3=\nabla_{E_3}E_2=\tfrac12E_1,
\qquad \nabla_{E_i}E_i=0.
\end{aligned}
\]
Przy konwencji $R(X,Y)Z=\nabla_X\nabla_YZ-\nabla_Y\nabla_XZ-
\nabla_{[X,Y]}Z$ otrzymujemy na przykład
$R(E_1,E_2)E_2=-\tfrac34E_1$ oraz
$R(E_1,E_3)E_3=\tfrac14E_1$. Trzeci rachunek jest analogiczny.
Tensor Ricciego jest sumą krzywizn przekrojowych po ortonormalnych
kierunkach; daje to podane wartości.
\end{proof}

\begin{przyklad}[Zwarta nilrozmaitość]
Podgrupa $\Gamma=H_3(\mathbb Z)$ macierzy Heisenberga o całkowitych
współrzędnych jest dyskretna i współzwarta. Iloraz
$\Gamma\backslash H_3(\mathbb R)$ jest zwartą trójrozmaitością
z geometrią Nil. Stanowi nietrywialną wiązkę okręgów nad
$\mathbb T^2$: przechodząc wokół podstawowego równoległoboku w bazie,
włókno doznaje przesunięcia wynikającego z niezerowej krzywizny
połączenia. Nie jest to torus płaski $\mathbb T^3$.
\end{przyklad}

\section{Geometria \texorpdfstring{$\widetilde{\mathrm{SL}}_2(\mathbb R)$}{SL2-tylda}}
\label{sec:thurston-sl2}

Geometria ta jest naturalnie związana z jednostkową wiązką styczną
płaszczyzny hiperbolicznej. Grupa izometrii zachowujących orientację
$\mathbb H^2$ jest izomorficzna z
$\operatorname{PSL}_2(\mathbb R)$; działa ona swobodnie i tranzytywnie
na wektorach stycznych długości jeden. Stąd
$T^1\mathbb H^2\simeq\operatorname{PSL}_2(\mathbb R)$ jako rozmaitości.
Nakryciem uniwersalnym tej przestrzeni jest
$\widetilde{\operatorname{SL}}_2(\mathbb R)$.

W modelu górnej półpłaszczyzny $\mathbb H^2=\{(x,y):y>0\}$
można, po identyfikacji nakrycia uniwersalnego z
$\mathbb H^2\times\mathbb R$, wybrać metrykę
\begin{equation}\label{eq:sl2-metric}
 g_{\widetilde{\mathrm{SL}}_2}
 =\frac{dx^2+dy^2}{y^2}
 +\left(dt+\frac{dx}{y}\right)^2.
\end{equation}
Rzut na $(x,y)$ jest submersją Riemannowską o jednowymiarowych,
geodezyjnych włóknach. Forma połączenia
$\alpha=dt+dx/y$ ma krzywiznę
$d\alpha=dx\wedge dy/y^2$, a więc jest ona (z uwzględnieniem
orientacji) równa formie pola powierzchni hiperbolicznej.
Nie jest to metryka produktowa.

\begin{uwaga}[Krzywizny]
Dla normalizacji~\eqref{eq:sl2-metric} wzory O'Neilla dla
submersji Riemannowskiej dają krzywiznę poziomej płaszczyzny
$-1-3/4=-7/4$ oraz krzywizny płaszczyzn zawierających kierunek
włókna równe $1/4$. W szczególności w poziomej ortonormalnej
bazie $E_1,E_2$ i pionowym $E_3$ tensor Ricciego ma wartości
$(-3/2,-3/2,1/2)$. Konkretne wartości zależą od normalizacji
współczynnika przy $\alpha^2$; znak niezerowej krzywizny
połączenia jest ważniejszy od tych liczb.
\end{uwaga}

\begin{przyklad}[Wiązka jednostkowych wektorów nad powierzchnią]
Niech $\Sigma_g=\Gamma\backslash\mathbb H^2$ będzie zamkniętą,
zorientowaną powierzchnią hiperboliczną rodzaju $g\ge2$.
Wówczas $T^1\Sigma_g$ jest zwartą trójrozmaitością, która otrzymuje
geometrię $\widetilde{\operatorname{SL}}_2(\mathbb R)$.
Wiązka $\mathbb S^1\to T^1\Sigma_g\to\Sigma_g$ ma liczbę
Eulera $\chi(\Sigma_g)=2-2g\neq0$. Wynika to z faktu, że
$T^1\Sigma_g$ jest wiązką okręgów stowarzyszoną ze zorientowaną
wiązką styczną $T\Sigma_g$, której klasa Eulera paruje z klasą
fundamentalną do charakterystyki Eulera powierzchni.
\end{przyklad}

Powyższy przykład należy porównać z
$\Sigma_g\times\mathbb S^1$ o geometrii
$\mathbb H^2\times\mathbb R$. Obie rozmaitości mają wiązkę okręgów
nad tą samą powierzchnią, ale w pierwszej liczba Eulera wiązki
jest niezerowa, zaś w drugim przypadku zerowa. Topologia
wiązki wykrywa więc różnicę między dwiema geometriami.

\section{Geometria Sol i odwzorowania sklejenia torusa}
\label{sec:thurston-sol}

Grupę Sol można zrealizować jako $\mathbb R^2\rtimes\mathbb R$,
gdzie $t\in\mathbb R$ działa na $\mathbb R^2$ macierzą
$\operatorname{diag}(e^t,e^{-t})$. Prawo grupowe jest dane wzorem
\[
(x,y,t)\cdot(x',y',t')
=(x+e^t x',\,y+e^{-t}y',\,t+t').
\]
W odróżnieniu od Nil, ta grupa jest rozwiązalna, ale nie nilpotentna.
Jednoformy $\theta^1=e^{-t}dx$, $\theta^2=e^t dy$ i
$\theta^3=dt$ są lewostronnie niezmiennicze. Wybieramy metrykę
\begin{equation}\label{eq:sol-metric}
 g_{\mathrm{Sol}}=e^{-2t}dx^2+e^{2t}dy^2+dt^2.
\end{equation}
Jej globalna ortonormalna rama to
$E_1=e^t\partial_x$, $E_2=e^{-t}\partial_y$,
$E_3=\partial_t$; spełnia ona
$[E_3,E_1]=E_1$, $[E_3,E_2]=-E_2$ oraz $[E_1,E_2]=0$.

\begin{twierdzenie}[Krzywizny geometrii Sol]
Dla metryki~\eqref{eq:sol-metric} zachodzi
$K_{12}=1$ oraz $K_{13}=K_{23}=-1$.
W szczególności
$\operatorname{Ric}=\operatorname{diag}(0,0,-2)$
w ramie $E_1,E_2,E_3$ i $\operatorname{Scal}=-2$.
\end{twierdzenie}
\begin{proof}
Ze wzoru Koszula otrzymujemy
$\nabla_{E_1}E_1=E_3$,
$\nabla_{E_1}E_3=-E_1$,
$\nabla_{E_2}E_2=-E_3$,
$\nabla_{E_2}E_3=E_2$;
pozostałe pochodne kowariantne pól ramy są zerowe.
Stąd $R(E_1,E_2)E_2=E_1$ oraz
$R(E_1,E_3)E_3=-E_1$ i $R(E_2,E_3)E_3=-E_2$.
Pozostałe wzory wynikają z definicji krzywizny Ricciego.
\end{proof}

\begin{przyklad}[Zwarta geometria Sol jako torus odwzorowania]
\label{ex:sol-mapping-torus}
Niech
$A=\left(\begin{smallmatrix}2&1\\1&1\end{smallmatrix}\right)
\in\mathrm{SL}_2(\mathbb Z)$. Jej wartości własne są rzeczywiste,
dodatnie i równe
$\lambda=(3+\sqrt5)/2>1$ oraz $\lambda^{-1}$.
Rozważmy \emph{torus odwzorowania}
\[
M_A=\frac{\mathbb T^2\times[0,1]}
{(u,1)\sim(Au,0)}.
\]
Jest to zwarta trójrozmaitość będąca wiązką torusów nad
$\mathbb S^1$. Aby zobaczyć geometrię Sol, wybierzmy liniowe
współrzędne na $\mathbb R^2$ diagonalizujące $A$.
W tych współrzędnych działanie $A$ ma postać
$(x,y)\mapsto(\lambda x,\lambda^{-1}y)$, natomiast
$\mathbb Z^2$ przechodzi w pewną kratę $\Lambda\subset\mathbb R^2$,
niezmienniczą względem tego przekształcenia.
Grupa generowana przez translacje o wektory $\Lambda$ oraz
\[
 (x,y,t)\longmapsto
 (\lambda x,\lambda^{-1}y,t+\log\lambda)
\]
działa na $(\mathrm{Sol},g_{\mathrm{Sol}})$ izometrycznie,
swobodnie i współzwarto. Otrzymany iloraz jest torusem
odwzorowania z monodromią $A$ (ewentualnie $A^{-1}$,
zależnie od konwencji sklejania końców przedziału).
Zatem $M_A$ ma geometrię Sol.
\end{przyklad}

Widać tu bezpośrednio związek algebry z geometrią: hiperboliczne
odwzorowanie torusa ma dwa kierunki własne, z których jeden ulega
rozciąganiu, a drugi ściskaniu. Metryka Sol dokładnie kompensuje te
zmiany przy przesunięciu w kierunku trzeciej współrzędnej.

\section{Wiązki Seiferta a geometrie modelowe}

\begin{definicja}[Wiązka Seiferta]
W uproszczeniu \emph{rozmaitością włóknistą Seiferta} nazywamy
trójrozmaitość rozłożoną na okręgi tak, że każde włókno ma
otoczenie modelowane na odpowiednim włóknistym torusie pełnym.
Włókna wyjątkowe dopuszczają skończone lokalne stabilizatory;
przestrzeń ilorazowa jest wówczas dwuwymiarowym orbifoldem.
Gdy włókien wyjątkowych brak, mamy zwykłą wiązkę okręgów nad
powierzchnią.
\end{definicja}

Przykłady $\mathbb S^2\times\mathbb S^1$,
$\Sigma_g\times\mathbb S^1$, zwarta nilrozmaitość oraz
$T^1\Sigma_g$ pokazują wspólny wzorzec: większość geometrii
Thurstona z jednowymiarową izotropią występuje naturalnie na
rozmaitościach z włóknowaniem przez okręgi.
Odpowiednia geometria zależy od krzywizny orbifoldu bazy
oraz liczby Eulera włóknowania. Dla zorientowanej wiązki
okręgów nad zamkniętą zorientowaną powierzchnią sytuację
ilustruje poniższe zestawienie; w ogólności trzeba uwzględnić
również włókna wyjątkowe.

\begin{center}
\begin{tabular}{|c|c|l|}
\hline
Krzywizna bazy & Liczba Eulera wiązki & Geometria \\
\hline
$+1$ & $0$ & $\mathbb S^2\times\mathbb R$ \\
$0$ & $0$ & $\mathbb R^3$ \\
$0$ & $\neq0$ & Nil \\
$-1$ & $0$ & $\mathbb H^2\times\mathbb R$ \\
$-1$ & $\neq0$ & $\widetilde{\mathrm{SL}}_2(\mathbb R)$ \\
\hline
\end{tabular}
\end{center}

Tabela nie obejmuje sferycznych przestrzeni z niezerową klasą
Eulera, które również bywają włókniste Seiferta. Pełne kryteria
dla orbifoldów wykorzystują ich orbifoldową charakterystykę
Eulera oraz odpowiednio znormalizowaną liczbę Eulera
włóknowania. W wersji ogólnej dla zamkniętych, orientowalnych
rozmaitości Seiferta przy $\chi^{\mathrm{orb}}>0$ występuje
geometria $\mathbb S^3$ lub $\mathbb S^2\times\mathbb R$;
przy $\chi^{\mathrm{orb}}=0$ geometria Nil lub $\mathbb R^3$;
przy $\chi^{\mathrm{orb}}<0$ geometria
$\widetilde{\mathrm{SL}}_2(\mathbb R)$ lub
$\mathbb H^2\times\mathbb R$, z rozróżnieniem według
zanikania liczby Eulera włóknowania.

\section{Od Ricciego do pełnej krzywizny w wymiarze trzecim}
\label{sec:thurston-ricci}

Wymiar trzy ma jeszcze jedną niezwykłą własność: tensor Riemanna
jest całkowicie określony przez tensor Ricciego. Jest to konsekwencja
zanikania tensora Weyla w wymiarze trzy.

\begin{twierdzenie}[Ricci determinuje krzywiznę w wymiarze trzy]
Niech $(M^3,g)$ będzie trójwymiarową rozmaitością Riemanna.
W każdym punkcie można wybrać ortonormalną bazę
$e_1,e_2,e_3$ diagonalizującą tensor Ricciego, o wartościach
własnych $r_1,r_2,r_3$. Wtedy krzywizny trzech płaszczyzn
własnych są równe
\[
K_{12}=\frac{r_1+r_2-r_3}{2},\quad
K_{13}=\frac{r_1+r_3-r_2}{2},\quad
K_{23}=\frac{r_2+r_3-r_1}{2}.
\]
W szczególności $\operatorname{Ric}=0$ wtedy i tylko wtedy,
gdy pełny tensor krzywizny Riemanna jest zerowy.
\end{twierdzenie}
\begin{proof}
Przypomnijmy, że
$r_1=K_{12}+K_{13}$,
$r_2=K_{12}+K_{23}$ i
$r_3=K_{13}+K_{23}$.
Rozwiązanie tego układu liniowego daje podane wzory.
W wymiarze trzy wszystkie składowe tensora krzywizny można
odzyskać z tensora Ricciego (równoważnie: część Weyla jest
identycznie zerowa). W szczególności zerowanie się wszystkich
wartości własnych tensora Ricciego daje zerowanie się tensora
Riemanna.
\end{proof}

\begin{wniosek}
Jeżeli trójwymiarowa rozmaitość Riemanna spełnia
$\operatorname{Ric}=\lambda g$ dla pewnej stałej $\lambda$,
to ma stałą krzywiznę przekrojową $K=\lambda/2$.
Każda zupełna jednospójna trójrozmaitość spełniająca
$\operatorname{Ric}=0$ jest izometryczna z $\mathbb R^3$.
\end{wniosek}
\begin{proof}
Pierwsza teza wynika bezpośrednio z powyższych wzorów.
W drugim przypadku metryka jest płaska; zastosowanie
klasyfikacji zupełnych jednospójnych przestrzeni o stałej
krzywiźnie daje $\mathbb R^3$.
\end{proof}

\begin{uwaga}[Co zmienia się w wyższych wymiarach?]
Dla $n\ge4$ tensor Weyla może być niezerowy, więc metryka
Ricci-płaska nie musi być płaska. Własność wykorzystana
powyżej jest specyficzna dla wymiaru trzy i wyjaśnia,
dlaczego warunek $\operatorname{Ric}=0$ ma w geometrii
trójwymiarowej znacznie silniejsze konsekwencje niż w
geometrii czterowymiarowej.
\end{uwaga}

\section{Hiperboliczne trójrozmaitości i wielościany fundamentalne}
\label{sec:thurston-hyperbolic}

Konstrukcja zwartych ilorazów $\mathbb H^3$ wymaga innej techniki
niż torus odwzorowania czy kratowe ilorazy grup Liego.
Naturalnym narzędziem są wielościany fundamentalne i izometrie
parujące ich ściany.

\begin{definicja}[Dziedzina fundamentalna]
Niech $\Gamma$ będzie dyskretną grupą izometrii $X$.
Domknięty zbiór $P\subset X$ nazywamy \emph{dziedziną
fundamentalną}, jeżeli $\Gamma P=X$, a wnętrza jego różnych
translat są rozłączne, tzn.
$\operatorname{int}(P)\cap\gamma\operatorname{int}(P)=\varnothing$
dla $\gamma\in\Gamma\setminus\{e\}$.
\end{definicja}

Twierdzenie Poincarégo o wielościanach daje warunki, pod którymi
przyporządkowanie izometrii parom ścian wielościanu definiuje
dyskretną grupę, a sam wielościan jest fundamentalny. W pełnej
wersji należy zweryfikować zgodność parowania ścian,
relacje wokół krawędzi i warunki lokalnego domknięcia oraz
zupełności. Sam fakt, że ściany dają się sparować, nie
wystarcza do otrzymania poprawnej rozmaitości hiperbolicznej.

\begin{przyklad}[Dopełnienie węzła ósemkowego]
Niech $K\subset\mathbb S^3$ będzie węzłem ósemkowym.
Wnętrze dopełnienia $\mathbb S^3\setminus N(K)$, gdzie $N(K)$
jest regularnym otoczeniem węzła, ma zupełną strukturę
hiperboliczną o skończonej objętości. Można ją zbudować przez
sklejenie dwóch regularnych idealnych czworościanów
hiperbolicznych. Jeśli $v_3$ oznacza objętość regularnego
idealnego czworościanu, to objętość tej rozmaitości wynosi
$2v_3$. Nie jest to przykład \emph{zwartej} rozmaitości:
na końcu znajduje się cusp (koniec paraboliczny) z
przekrojem będącym dwuwymiarowym torusem.
\end{przyklad}

\begin{przyklad}[Rozmaitość Giesekinga]
Rozmaitość Giesekinga można otrzymać przez odpowiednie
identyfikacje ścian jednego regularnego idealnego
czworościanu hiperbolicznego. Jest ona nieorientowalna,
niezwarta, zupełna i ma skończoną objętość $v_3$.
Jej orientowalne nakrycie dwukrotne jest izometryczne z
dopełnieniem węzła ósemkowego (z naturalną metryką
hiperboliczną). Ten przykład pokazuje, że topologia końców
i orientowalność są istotne nawet przy stałej krzywiźnie $-1$.
\end{przyklad}

\section{Klasyfikacja lokalnych geometrii a geometryzacja}

Nie należy utożsamiać dwóch twierdzeń: klasyfikacja Thurstona
opisuje osiem możliwych \emph{maksymalnych geometrii modelowych}
na jednospójnych jednorodnych przestrzeniach, natomiast
\emph{geometryzacja} dotyczy globalnej struktury topologicznej
dowolnych trójrozmaitości.

\begin{twierdzenie}[Geometryzacja Thurstona--Perelmana, schemat]
Każda zwarta orientowalna trójrozmaitość po rozkładzie na
składniki pierwsze (cięcia wzdłuż sfer dwuwymiarowych),
a następnie odpowiednim rozkładzie wzdłuż nieściśliwych
torusów, daje części typu Seiferta i części hiperboliczne.
Hiperboliczne wnętrza odpowiednich części mają zupełne
metryki hiperboliczne o skończonej objętości.
Wśród możliwych geometrii składników zamkniętych i
przykładów podstawowych pojawia się osiem geometrii
z twierdzenia~\ref{thm:osiem-geometrii}.
\end{twierdzenie}

Sformułowanie jest celowo schematyczne: poprawna wersja
wymaga uwzględnienia redukcji sferycznej, rozkładu JSJ,
przypadków o pustym rozkładzie i sposobu traktowania brzegu.
Dowód geometryzacji nie jest konsekwencją samej listy ośmiu
geometrii. Historycznie przełomem było zastosowanie przepływu
Ricciego z chirurgią; schemat rozwoju metryki zapisuje
równanie
\[
 \frac{\partial g(t)}{\partial t}=-2\operatorname{Ric}(g(t)).
\]
Jego szczegółowa analiza wymaga oddzielnego rozdziału.

\begin{uwaga}[Różne znaczenia słowa „geometria”]
Nie każda zwarta trójrozmaitość dopuszcza jedną geometrię
jednorodną na całej swojej przestrzeni. Geometryzacja mówi,
że po odpowiednich cięciach otrzymujemy części modelowane
przez geometrie z listy. Jest to zasadnicza różnica między
rozmaitością \emph{geometryczną} a rozmaitością
\emph{geometryzowalną}.
\end{uwaga}

\section{Podsumowanie i powiązania z poprzednimi rozdziałami}

Osiem geometrii Thurstona powstaje z dwóch współdziałających
idei: jednorodności metrycznej i istnienia zwartych ilorazów.
Z punktu widzenia rachunku różniczkowego najważniejsze są
globalne ortonormalne ramy, wzór Koszula i obliczenia krzywizn
Nil oraz Sol. Z punktu widzenia topologii zasadniczą rolę
odgrywają grupy podstawowe, nakrycia, wiązki okręgów, liczba
Eulera i torusy odwzorowania. Geometria
$\widetilde{\mathrm{SL}}_2(\mathbb R)$ łączy oba punkty widzenia
przez krzywiznę połączenia wiązki Seiferta.

Warto przy tym pamiętać, że twierdzenie o geometryzacji
ma charakter globalny i nie wynika z lokalnej listy geometrii.
Jest ono pomostem do bardziej zaawansowanej topologii
trójrozmaitości, teorii rozkładów JSJ i przepływu Ricciego.

\paragraph{Literatura uzupełniająca.}
W.~P.~Thurston, \emph{Three-Dimensional Geometry and Topology},
vol.~1, Princeton University Press, 1997;
P.~Scott, \emph{The Geometries of 3-Manifolds},
Bulletin of the London Mathematical Society 15 (1983), 401--487;
J.~W.~Milnor, \emph{Curvatures of Left Invariant Metrics on Lie Groups},
Advances in Mathematics 21 (1976), 293--329;
S.~A.~Pinillos Prado, \emph{Geometric Structures on 3-Dimensional
Manifolds}, Universidade Federal de Minas Gerais, 2021.

% Wymagane w preambule: amsmath, amssymb, amsthm, mathtools.
% Srodowiska: definicja, twierdzenie, przyklad, uwaga, wniosek, proof.
% Jezeli w skrypcie stosowana jest inna definicja srodowisk,
% nalezy dostosowac nazwy w preambule, nie w tresci rozdzialu.

\chapter{Geometrie Thurstona i struktury geometryczne na trójrozmaitościach}
\label{chap:geometrie-thurstona}

W poprzednich rozdziałach rozważaliśmy rozmaitości Riemanna, ich połączenia,
geodezyjne i krzywiznę oraz podstawowe geometrie przestrzeni o stałej krzywiźnie.
Tutaj przejdziemy od geometrii pojedynczej rozmaitości do pytania o to,
\emph{jakie geometrie jednorodne mogą występować lokalnie na trójrozmaitościach}.
Z jednej strony prowadzi to do klasyfikacji ośmiu geometrii modelowych
Thurstona, z drugiej -- do związków między metryką a topologią rozmaitości
trójwymiarowych. W szczególności zobaczymy, dlaczego do opisu wszystkich
zwartych trójrozmaitości nie wystarczają geometrie sferyczna, euklidesowa
i hiperboliczna.

\section{Struktury \texorpdfstring{$(G,X)$}{(G,X)} i geometrie modelowe}

Ideą programu erlangeńskiego Kleina jest badanie przestrzeni za pomocą
grupy przekształceń zachowujących interesującą nas strukturę. Geometria
Riemanna dostarcza szczególnego przypadku: jeśli $(X,g)$ jest rozmaitością
Riemanna, to grupa $\operatorname{Isom}(X,g)$ składa się z jej izometrii.

\begin{definicja}[Struktura $(G,X)$]
Niech $X$ będzie spójną gładką rozmaitością, a $G$ grupą Liego działającą na
$X$ przez dyfeomorfizmy. \emph{Strukturą $(G,X)$} na rozmaitości $M$
nazywamy atlas map $\{(U_\alpha,\varphi_\alpha)\}$,
gdzie $\varphi_\alpha\colon U_\alpha\to V_\alpha\subset X$ jest
dyfeomorfizmem na zbiór otwarty, a dla każdej spójnej składowej
$C\subset U_\alpha\cap U_\beta$ istnieje $g_{\alpha\beta,C}\in G$ takie, że
$\varphi_\alpha\circ\varphi_\beta^{-1}$ na $\varphi_\beta(C)$
jest obcięciem działania $g_{\alpha\beta,C}$.
\end{definicja}

Warunek dotyczący składowych przecięcia jest istotny: jedno przecięcie
$U_\alpha\cap U_\beta$ może być niespójne, a na różnych jego składowych
mogą występować różne elementy grupy $G$. Zwykły atlas gładki odpowiada
pseudogrupie wszystkich lokalnych dyfeomorfizmów, podczas gdy struktura
$(G,X)$ dopuszcza tylko lokalne przekształcenia pochodzące z jednej grupy.

\begin{przyklad}[Torus płaski]
Niech $X=\mathbb R^3$, $G=\operatorname{Isom}(\mathbb R^3)$ i
$\Gamma=\mathbb Z^3$ działająca przez translacje. Iloraz
$\mathbb T^3=\mathbb R^3/\mathbb Z^3$ otrzymuje atlas $(G,X)$, ponieważ
lokalne odwrotności nakrycia różnią się o translacje należące do $G$.
Wszystkie przestrzenie styczne są euklidesowe, chociaż $\mathbb T^3$
nie jest globalnie izometryczny z $\mathbb R^3$.
\end{przyklad}

\begin{definicja}[Geometria modelowa Thurstona]
\label{def:model-thurston}
Przez \emph{geometrię modelową} rozumiemy parę $(G,X)$, dla której:
\begin{enumerate}
\item $X$ jest spójna i jednospójna;
\item $G$ działa na $X$ tranzytywnie, a stabilizator punktu
      $G_x=\{a\in G:a\cdot x=x\}$ jest zwarty;
\item $G$ jest maksymalną grupą dyfeomorfizmów $X$ spełniającą warunek
      zwartości stabilizatorów;
\item istnieje zwarta rozmaitość z pełną strukturą $(G,X)$.
\end{enumerate}
Dwie geometrie utożsamiamy, gdy ich przestrzenie i grupy działań są
izomorficzne w naturalnym sensie. Przymiotnik \emph{pełna} oznacza, że
struktura pochodzi z ilorazu $X/\Gamma$ dla odpowiednio działającej
grupy dyskretnej $\Gamma\subset G$.
\end{definicja}

Dwa ostatnie warunki nie są kosmetyczne. Gdyby zrezygnować z maksymalności,
tę samą geometrię euklidesową można by sztucznie traktować jako wiele
różnych geometrii, wybierając mniejsze grupy ruchów. Wymóg zwartego
ilorazu eliminuje z kolei liczne jednorodne przestrzenie niemające
zwartych realizacji globalnych.

\begin{twierdzenie}[Metryka niezmiennicza]
\label{thm:metryka-niezmiennicza}
Jeżeli grupa Liego $G$ działa tranzytywnie na $X$, a stabilizator $G_o$
pewnego punktu $o\in X$ jest zwarty, to na $X$ istnieje metryka
Riemanna niezmiennicza względem $G$.
\end{twierdzenie}
\begin{proof}
Wybierzmy dowolny iloczyn skalarny $h$ na $T_oX$. Niech $\mu$ będzie
unormowaną miarą Haara na zwartej grupie $G_o$. Definiujemy
\[
 h_o^{G}(u,v)=\int_{G_o}h(da_o u,da_o v)\,d\mu(a).
\]
Otrzymana forma jest dodatnio określona i niezmiennicza względem
reprezentacji izotropii. Jeśli $x=b\cdot o$, przenosimy ją wzorem
$h_x^G(db_o u,db_o v)=h_o^G(u,v)$. Definicja nie zależy od wyboru $b$:
dwa takie elementy różnią się o element stabilizatora $G_o$, a forma
$h_o^G$ została względem niego uśredniona. Gładkość metryki wynika z
gładkości działania i lokalnego wyboru przekrojów odwzorowania
$G\to G/G_o\simeq X$.
\end{proof}

\begin{uwaga}[Nakrycie uniwersalne i grupa decka]
Jeżeli $X$ jest jednospójna, $\Gamma\subset\operatorname{Isom}(X,g)$
działa swobodnie, właściwie nieciągle i współzwarto, to
$M=\Gamma\backslash X$ jest zwartą rozmaitością Riemanna, a
$X\to M$ jest jej nakryciem uniwersalnym. W szczególności
$\pi_1(M)\simeq\Gamma$. Odwrotnie, każda \emph{pełna} struktura
$(G,X)$ na zwartej rozmaitości jest realizowana w ten sposób.
Nie każda struktura $(G,X)$ jest jednak z definicji pełna.
\end{uwaga}

Przydatne pojęcie pośrednie daje \emph{odwzorowanie rozwijające}
$\operatorname{dev}\colon\widetilde M\to X$. Powstaje ono przez analityczne
kontynuowanie kolejnych lokalnych map modelowych na nakryciu
uniwersalnym i spełnia
$\operatorname{dev}(\gamma p)=\rho(\gamma)\operatorname{dev}(p)$,
gdzie $\rho\colon\pi_1(M)\to G$ jest \emph{holonomią}. Jeżeli struktura
jest pełna, odwzorowanie rozwijające identyfikuje $\widetilde M$ z $X$.

\section{Dlaczego w wymiarze trzecim pojawia się osiem geometrii?}

\begin{twierdzenie}[Klasyfikacja geometrii modelowych Thurstona]
\label{thm:osiem-geometrii}
W wymiarze trzecim, z dokładnością do izomorfizmu geometrii modelowych,
istnieje dokładnie osiem możliwości:
\[
\mathbb S^3,\qquad \mathbb R^3,\qquad \mathbb H^3,\qquad
\mathbb S^2\!\times\!\mathbb R,\qquad
\mathbb H^2\!\times\!\mathbb R,\qquad
\mathrm{Nil},\qquad\widetilde{\mathrm{SL}}_2(\mathbb R),\qquad
\mathrm{Sol}.
\]
Dla trzech pierwszych geometrii spójna izotropia jest trójwymiarowa;
dla następnych czterech -- jednowymiarowa; dla geometrii Sol --
zerowymiarowa.
\end{twierdzenie}

Podamy mechanizm klasyfikacji, oddzielając go od pełnego dowodu, który
wymaga dodatkowej analizy jednorodnych grup Liego i ich sieci.
Na mocy twierdzenia~\ref{thm:metryka-niezmiennicza} możemy wybrać
metrykę niezmienniczą na $X$. Izotropia działa wiernie na $T_oX$
(izometria ustalająca punkt i mająca identyczną różniczkę jest
identycznością), zatem jej składowa spójna jest zwartą spójną
podgrupą $\mathrm{SO}(3)$. Takie podgrupy są, z dokładnością do
sprzężenia, równe $\mathrm{SO}(3)$, $\mathrm{SO}(2)$ albo grupie
trywialnej.

Gdy izotropia ma wymiar trzy, działa tranzytywnie na kierunkach
w każdej przestrzeni stycznej. Krzywizna przekrojowa zależy więc
tylko od punktu, a jednorodność czyni ją stałą na całym $X$.
Otrzymujemy trzy przestrzenie stałej krzywizny.

Gdy izotropia ma wymiar jeden, zachowuje pewien wyróżniony kierunek
pola prostych. Na odpowiednim nakryciu pole to można zorientować i
otrzymać niezmiennicze pole wektorowe, którego linie całkowe tworzą
jednowymiarowe włóknowanie. Geometria poprzeczna jest jedną z trzech
geometrii powierzchni jednorodnych; krzywizna połączenia tego
włóknowania odróżnia geometrie produktowe od Nil i
$\widetilde{\mathrm{SL}}_2(\mathbb R)$. Przy zerowymiarowej izotropii
analizuje się trójwymiarowe grupy Liego wyposażone w metryki
lewostronnie niezmiennicze. Po uwzględnieniu maksymalności grupy
izometrii i istnienia zwartego ilorazu jedyną nową geometrią jest Sol.
To jest \emph{szkic dowodu}, nie pełna klasyfikacja możliwych algebr
Liego i ich sieci.

\begin{center}
\begin{tabular}{|l|c|l|}
\hline
Geometria & $\dim G_o^0$ & Typ geometrii poprzecznej lub krzywizna \\
\hline
$\mathbb S^3$ & $3$ & $K\equiv 1$ \\
$\mathbb R^3$ & $3$ & $K\equiv 0$ \\
$\mathbb H^3$ & $3$ & $K\equiv -1$ \\
$\mathbb S^2\times\mathbb R$ & $1$ & produkt, baza $\mathbb S^2$ \\
$\mathbb H^2\times\mathbb R$ & $1$ & produkt, baza $\mathbb H^2$ \\
Nil & $1$ & baza $\mathbb R^2$, skręcone włókna \\
$\widetilde{\mathrm{SL}}_2(\mathbb R)$ & $1$ & baza $\mathbb H^2$, skręcone włókna \\
Sol & $0$ & włóknowanie nad $\mathbb R$ \\
\hline
\end{tabular}
\end{center}

\section{Trzy geometrie o stałej krzywiźnie}

Przypomnijmy jedynie elementy potrzebne do dalszego porównania. Po
odpowiednim przeskalowaniu metryk geometrie te mają krzywiznę
przekrojową odpowiednio $\kappa=1,0,-1$. W wymiarze trzy ich
tensor Ricciego i krzywizna skalarna spełniają
$\operatorname{Ric}=2\kappa g$ oraz $\operatorname{Scal}=6\kappa$.

\begin{przyklad}[Modele i zwarte ilorazy]
\begin{enumerate}
\item $\mathbb S^3=\{z\in\mathbb C^2:|z_1|^2+|z_2|^2=1\}$,
      z metryką indukowaną z $\mathbb R^4$. Przykładami zwartych
      ilorazów są przestrzenie soczewkowe
      $L(p,q)=\mathbb S^3/\mathbb Z_p$, gdzie generator działa przez
      $(z_1,z_2)\mapsto(e^{2\pi i/p}z_1,e^{2\pi iq/p}z_2)$,
      a $\gcd(p,q)=1$.
\item $\mathbb R^3$, z metryką $g=dx^2+dy^2+dz^2$;
      ilorazem zwartym jest $\mathbb T^3$.
\item $\mathbb H^3=\{(x,y,z)\in\mathbb R^3:z>0\}$,
      z metryką $g=z^{-2}(dx^2+dy^2+dz^2)$. Istnieją zwarte ilorazy
      $\Gamma\backslash\mathbb H^3$, gdzie $\Gamma$ jest bezskrętną
      współzwartą kratą w $\operatorname{Isom}(\mathbb H^3)$.
\end{enumerate}
\end{przyklad}

\section{Dwie geometrie produktowe}

Niech $(\Sigma^2,g_\Sigma)$ będzie powierzchnią o stałej krzywiźnie
Gaussa $\kappa\in\{1,-1\}$. Na $X=\Sigma\times\mathbb R$ bierzemy
metrykę produktową $g=g_\Sigma+dt^2$. Ponieważ połączenie
Levi-Civity produktu zachowuje oba czynniki, dla ortonormalnych
wektorów $e_1,e_2$ stycznych do $\Sigma$ oraz $e_3=\partial_t$
mamy $K(e_1,e_2)=\kappa$ i $K(e_1,e_3)=K(e_2,e_3)=0$.
W konsekwencji
\[
\operatorname{Ric}(e_1,e_1)=\operatorname{Ric}(e_2,e_2)=\kappa,
\qquad \operatorname{Ric}(e_3,e_3)=0,
\qquad \operatorname{Scal}=2\kappa.
\]
Są to geometrie jednorodne, lecz \emph{nie} mają stałej krzywizny
przekrojowej w wymiarze trzy.

\begin{przyklad}[Zwarte realizacje]
Dla $\mathbb S^2\times\mathbb R$ otrzymujemy
$\mathbb S^2\times\mathbb S^1$ przez ilorazowanie drugiego czynnika
translacjami całkowitymi. Dla $\mathbb H^2\times\mathbb R$ wybieramy
zwartą powierzchnię hiperboliczną $\Sigma_g=\Gamma\backslash\mathbb H^2$
rodzaju $g\ge2$; wtedy $\Sigma_g\times\mathbb S^1$ jest zwartym
ilorazem $\mathbb H^2\times\mathbb R$.
\end{przyklad}

\section{Geometria Nil: grupa Heisenberga}
\label{sec:thurston-nil}

Pierwsza istotnie nowa geometria pojawia się wtedy, gdy płaszczyznę
$\mathbb R^2$ zastąpimy rodziną prostych, których połączenie ma
niezerową krzywiznę. Rozważmy $\mathbb R^3$ z działaniem grupowym
\[
(x,y,z)\cdot(x',y',z')=
(x+x',\,y+y',\,z+z'+xy').
\]
Jest to trójwymiarowa \emph{grupa Heisenberga} $H_3$. Można ją
utożsamić z grupą macierzy
$\left(\begin{smallmatrix}1&x&z\\0&1&y\\0&0&1\end{smallmatrix}\right)$.

Jednoformy $\theta^1=dx$, $\theta^2=dy$ i
$\theta^3=dz-x\,dy$ są lewostronnie niezmiennicze.
Ich formami dualnymi są pola
$E_1=\partial_x$, $E_2=\partial_y+x\partial_z$ oraz
$E_3=\partial_z$. Bezpośrednio
$[E_1,E_2]=E_3$, zaś pozostałe nawiasy są zerowe.
Nadajmy $H_3$ metrykę
\begin{equation}\label{eq:nil-metric}
 g_{\mathrm{Nil}}=(\theta^1)^2+(\theta^2)^2+(\theta^3)^2
 =dx^2+dy^2+(dz-x\,dy)^2.
\end{equation}
Zapis $\theta^2$ w powyższej formule oznacza iloczyn tensorowy
$\theta\otimes\theta$, a nie iloczyn zewnętrzny. Pola $E_1,E_2,E_3$
tworzą globalną ortonormalną ramę.

\begin{przyklad}[Pozioma krzywa i holonomia]
Rzut $\pi\colon H_3\to\mathbb R^2$, $\pi(x,y,z)=(x,y)$,
ma włókna równoległe do osi $z$. Rozkład poziomy określa
$\theta^3=0$, czyli dla krzywej $\gamma(s)=(x(s),y(s),z(s))$
warunek poziomości brzmi $z'(s)=x(s)y'(s)$. Po przejściu
wzdłuż zamkniętej pętli $c$ w płaszczyźnie podstawa zostaje
niezmieniona, ale współrzędna $z$ zmienia się o
\[
\Delta z=\oint_c x\,dy=\iint_D dx\wedge dy,
\]
gdzie $D$ jest zorientowanym obszarem ograniczonym przez $c$.
Oznacza to, że przesunięcie wzdłuż włókna mierzy pole zorientowane
ograniczone przez pętlę. Ponieważ $d\theta^3=-dx\wedge dy$,
rozkład poziomy nie jest całkowalny.
\end{przyklad}

\begin{twierdzenie}[Krzywizny geometrii Nil]
Dla metryki~\eqref{eq:nil-metric} krzywizny przekrojowe płaszczyzn
wyznaczonych przez ramę $E_1,E_2,E_3$ wynoszą
$K_{12}=-3/4$ i $K_{13}=K_{23}=1/4$.
Stąd $\operatorname{Ric}=\operatorname{diag}(-1/2,-1/2,1/2)$
w tej ramie i $\operatorname{Scal}=-1/2$.
\end{twierdzenie}
\begin{proof}
Dla ortonormalnych pól lewostronnie niezmienniczych wzór Koszula
przybiera postać
\[
2\langle\nabla_XY,Z\rangle
=\langle[X,Y],Z\rangle-\langle[Y,Z],X\rangle
+\langle[Z,X],Y\rangle.
\]
Z jedynego niezerowego nawiasu $[E_1,E_2]=E_3$ dostajemy
\[
\begin{aligned}
&\nabla_{E_1}E_2=\tfrac12E_3,\qquad
 \nabla_{E_2}E_1=-\tfrac12E_3,\\
&\nabla_{E_1}E_3=\nabla_{E_3}E_1=-\tfrac12E_2,\\
&\nabla_{E_2}E_3=\nabla_{E_3}E_2=\tfrac12E_1,
\qquad \nabla_{E_i}E_i=0.
\end{aligned}
\]
Przy konwencji $R(X,Y)Z=\nabla_X\nabla_YZ-\nabla_Y\nabla_XZ-
\nabla_{[X,Y]}Z$ otrzymujemy na przykład
$R(E_1,E_2)E_2=-\tfrac34E_1$ oraz
$R(E_1,E_3)E_3=\tfrac14E_1$. Trzeci rachunek jest analogiczny.
Tensor Ricciego jest sumą krzywizn przekrojowych po ortonormalnych
kierunkach; daje to podane wartości.
\end{proof}

\begin{przyklad}[Zwarta nilrozmaitość]
Podgrupa $\Gamma=H_3(\mathbb Z)$ macierzy Heisenberga o całkowitych
współrzędnych jest dyskretna i współzwarta. Iloraz
$\Gamma\backslash H_3(\mathbb R)$ jest zwartą trójrozmaitością
z geometrią Nil. Stanowi nietrywialną wiązkę okręgów nad
$\mathbb T^2$: przechodząc wokół podstawowego równoległoboku w bazie,
włókno doznaje przesunięcia wynikającego z niezerowej krzywizny
połączenia. Nie jest to torus płaski $\mathbb T^3$.
\end{przyklad}

\section{Geometria \texorpdfstring{$\widetilde{\mathrm{SL}}_2(\mathbb R)$}{SL2-tylda}}
\label{sec:thurston-sl2}

Geometria ta jest naturalnie związana z jednostkową wiązką styczną
płaszczyzny hiperbolicznej. Grupa izometrii zachowujących orientację
$\mathbb H^2$ jest izomorficzna z
$\operatorname{PSL}_2(\mathbb R)$; działa ona swobodnie i tranzytywnie
na wektorach stycznych długości jeden. Stąd
$T^1\mathbb H^2\simeq\operatorname{PSL}_2(\mathbb R)$ jako rozmaitości.
Nakryciem uniwersalnym tej przestrzeni jest
$\widetilde{\operatorname{SL}}_2(\mathbb R)$.

W modelu górnej półpłaszczyzny $\mathbb H^2=\{(x,y):y>0\}$
można, po identyfikacji nakrycia uniwersalnego z
$\mathbb H^2\times\mathbb R$, wybrać metrykę
\begin{equation}\label{eq:sl2-metric}
 g_{\widetilde{\mathrm{SL}}_2}
 =\frac{dx^2+dy^2}{y^2}
 +\left(dt+\frac{dx}{y}\right)^2.
\end{equation}
Rzut na $(x,y)$ jest submersją Riemannowską o jednowymiarowych,
geodezyjnych włóknach. Forma połączenia
$\alpha=dt+dx/y$ ma krzywiznę
$d\alpha=dx\wedge dy/y^2$, a więc jest ona (z uwzględnieniem
orientacji) równa formie pola powierzchni hiperbolicznej.
Nie jest to metryka produktowa.

\begin{uwaga}[Krzywizny]
Dla normalizacji~\eqref{eq:sl2-metric} wzory O'Neilla dla
submersji Riemannowskiej dają krzywiznę poziomej płaszczyzny
$-1-3/4=-7/4$ oraz krzywizny płaszczyzn zawierających kierunek
włókna równe $1/4$. W szczególności w poziomej ortonormalnej
bazie $E_1,E_2$ i pionowym $E_3$ tensor Ricciego ma wartości
$(-3/2,-3/2,1/2)$. Konkretne wartości zależą od normalizacji
współczynnika przy $\alpha^2$; znak niezerowej krzywizny
połączenia jest ważniejszy od tych liczb.
\end{uwaga}

\begin{przyklad}[Wiązka jednostkowych wektorów nad powierzchnią]
Niech $\Sigma_g=\Gamma\backslash\mathbb H^2$ będzie zamkniętą,
zorientowaną powierzchnią hiperboliczną rodzaju $g\ge2$.
Wówczas $T^1\Sigma_g$ jest zwartą trójrozmaitością, która otrzymuje
geometrię $\widetilde{\operatorname{SL}}_2(\mathbb R)$.
Wiązka $\mathbb S^1\to T^1\Sigma_g\to\Sigma_g$ ma liczbę
Eulera $\chi(\Sigma_g)=2-2g\neq0$. Wynika to z faktu, że
$T^1\Sigma_g$ jest wiązką okręgów stowarzyszoną ze zorientowaną
wiązką styczną $T\Sigma_g$, której klasa Eulera paruje z klasą
fundamentalną do charakterystyki Eulera powierzchni.
\end{przyklad}

Powyższy przykład należy porównać z
$\Sigma_g\times\mathbb S^1$ o geometrii
$\mathbb H^2\times\mathbb R$. Obie rozmaitości mają wiązkę okręgów
nad tą samą powierzchnią, ale w pierwszej liczba Eulera wiązki
jest niezerowa, zaś w drugim przypadku zerowa. Topologia
wiązki wykrywa więc różnicę między dwiema geometriami.

\section{Geometria Sol i odwzorowania sklejenia torusa}
\label{sec:thurston-sol}

Grupę Sol można zrealizować jako $\mathbb R^2\rtimes\mathbb R$,
gdzie $t\in\mathbb R$ działa na $\mathbb R^2$ macierzą
$\operatorname{diag}(e^t,e^{-t})$. Prawo grupowe jest dane wzorem
\[
(x,y,t)\cdot(x',y',t')
=(x+e^t x',\,y+e^{-t}y',\,t+t').
\]
W odróżnieniu od Nil, ta grupa jest rozwiązalna, ale nie nilpotentna.
Jednoformy $\theta^1=e^{-t}dx$, $\theta^2=e^t dy$ i
$\theta^3=dt$ są lewostronnie niezmiennicze. Wybieramy metrykę
\begin{equation}\label{eq:sol-metric}
 g_{\mathrm{Sol}}=e^{-2t}dx^2+e^{2t}dy^2+dt^2.
\end{equation}
Jej globalna ortonormalna rama to
$E_1=e^t\partial_x$, $E_2=e^{-t}\partial_y$,
$E_3=\partial_t$; spełnia ona
$[E_3,E_1]=E_1$, $[E_3,E_2]=-E_2$ oraz $[E_1,E_2]=0$.

\begin{twierdzenie}[Krzywizny geometrii Sol]
Dla metryki~\eqref{eq:sol-metric} zachodzi
$K_{12}=1$ oraz $K_{13}=K_{23}=-1$.
W szczególności
$\operatorname{Ric}=\operatorname{diag}(0,0,-2)$
w ramie $E_1,E_2,E_3$ i $\operatorname{Scal}=-2$.
\end{twierdzenie}
\begin{proof}
Ze wzoru Koszula otrzymujemy
$\nabla_{E_1}E_1=E_3$,
$\nabla_{E_1}E_3=-E_1$,
$\nabla_{E_2}E_2=-E_3$,
$\nabla_{E_2}E_3=E_2$;
pozostałe pochodne kowariantne pól ramy są zerowe.
Stąd $R(E_1,E_2)E_2=E_1$ oraz
$R(E_1,E_3)E_3=-E_1$ i $R(E_2,E_3)E_3=-E_2$.
Pozostałe wzory wynikają z definicji krzywizny Ricciego.
\end{proof}

\begin{przyklad}[Zwarta geometria Sol jako torus odwzorowania]
\label{ex:sol-mapping-torus}
Niech
$A=\left(\begin{smallmatrix}2&1\\1&1\end{smallmatrix}\right)
\in\mathrm{SL}_2(\mathbb Z)$. Jej wartości własne są rzeczywiste,
dodatnie i równe
$\lambda=(3+\sqrt5)/2>1$ oraz $\lambda^{-1}$.
Rozważmy \emph{torus odwzorowania}
\[
M_A=\frac{\mathbb T^2\times[0,1]}
{(u,1)\sim(Au,0)}.
\]
Jest to zwarta trójrozmaitość będąca wiązką torusów nad
$\mathbb S^1$. Aby zobaczyć geometrię Sol, wybierzmy liniowe
współrzędne na $\mathbb R^2$ diagonalizujące $A$.
W tych współrzędnych działanie $A$ ma postać
$(x,y)\mapsto(\lambda x,\lambda^{-1}y)$, natomiast
$\mathbb Z^2$ przechodzi w pewną kratę $\Lambda\subset\mathbb R^2$,
niezmienniczą względem tego przekształcenia.
Grupa generowana przez translacje o wektory $\Lambda$ oraz
\[
 (x,y,t)\longmapsto
 (\lambda x,\lambda^{-1}y,t+\log\lambda)
\]
działa na $(\mathrm{Sol},g_{\mathrm{Sol}})$ izometrycznie,
swobodnie i współzwarto. Otrzymany iloraz jest torusem
odwzorowania z monodromią $A$ (ewentualnie $A^{-1}$,
zależnie od konwencji sklejania końców przedziału).
Zatem $M_A$ ma geometrię Sol.
\end{przyklad}

Widać tu bezpośrednio związek algebry z geometrią: hiperboliczne
odwzorowanie torusa ma dwa kierunki własne, z których jeden ulega
rozciąganiu, a drugi ściskaniu. Metryka Sol dokładnie kompensuje te
zmiany przy przesunięciu w kierunku trzeciej współrzędnej.

\section{Wiązki Seiferta a geometrie modelowe}

\begin{definicja}[Wiązka Seiferta]
W uproszczeniu \emph{rozmaitością włóknistą Seiferta} nazywamy
trójrozmaitość rozłożoną na okręgi tak, że każde włókno ma
otoczenie modelowane na odpowiednim włóknistym torusie pełnym.
Włókna wyjątkowe dopuszczają skończone lokalne stabilizatory;
przestrzeń ilorazowa jest wówczas dwuwymiarowym orbifoldem.
Gdy włókien wyjątkowych brak, mamy zwykłą wiązkę okręgów nad
powierzchnią.
\end{definicja}

Przykłady $\mathbb S^2\times\mathbb S^1$,
$\Sigma_g\times\mathbb S^1$, zwarta nilrozmaitość oraz
$T^1\Sigma_g$ pokazują wspólny wzorzec: większość geometrii
Thurstona z jednowymiarową izotropią występuje naturalnie na
rozmaitościach z włóknowaniem przez okręgi.
Odpowiednia geometria zależy od krzywizny orbifoldu bazy
oraz liczby Eulera włóknowania. Dla zorientowanej wiązki
okręgów nad zamkniętą zorientowaną powierzchnią sytuację
ilustruje poniższe zestawienie; w ogólności trzeba uwzględnić
również włókna wyjątkowe.

\begin{center}
\begin{tabular}{|c|c|l|}
\hline
Krzywizna bazy & Liczba Eulera wiązki & Geometria \\
\hline
$+1$ & $0$ & $\mathbb S^2\times\mathbb R$ \\
$0$ & $0$ & $\mathbb R^3$ \\
$0$ & $\neq0$ & Nil \\
$-1$ & $0$ & $\mathbb H^2\times\mathbb R$ \\
$-1$ & $\neq0$ & $\widetilde{\mathrm{SL}}_2(\mathbb R)$ \\
\hline
\end{tabular}
\end{center}

Tabela nie obejmuje sferycznych przestrzeni z niezerową klasą
Eulera, które również bywają włókniste Seiferta. Pełne kryteria
dla orbifoldów wykorzystują ich orbifoldową charakterystykę
Eulera oraz odpowiednio znormalizowaną liczbę Eulera
włóknowania. W wersji ogólnej dla zamkniętych, orientowalnych
rozmaitości Seiferta przy $\chi^{\mathrm{orb}}>0$ występuje
geometria $\mathbb S^3$ lub $\mathbb S^2\times\mathbb R$;
przy $\chi^{\mathrm{orb}}=0$ geometria Nil lub $\mathbb R^3$;
przy $\chi^{\mathrm{orb}}<0$ geometria
$\widetilde{\mathrm{SL}}_2(\mathbb R)$ lub
$\mathbb H^2\times\mathbb R$, z rozróżnieniem według
zanikania liczby Eulera włóknowania.

\section{Od Ricciego do pełnej krzywizny w wymiarze trzecim}
\label{sec:thurston-ricci}

Wymiar trzy ma jeszcze jedną niezwykłą własność: tensor Riemanna
jest całkowicie określony przez tensor Ricciego. Jest to konsekwencja
zanikania tensora Weyla w wymiarze trzy.

\begin{twierdzenie}[Ricci determinuje krzywiznę w wymiarze trzy]
Niech $(M^3,g)$ będzie trójwymiarową rozmaitością Riemanna.
W każdym punkcie można wybrać ortonormalną bazę
$e_1,e_2,e_3$ diagonalizującą tensor Ricciego, o wartościach
własnych $r_1,r_2,r_3$. Wtedy krzywizny trzech płaszczyzn
własnych są równe
\[
K_{12}=\frac{r_1+r_2-r_3}{2},\quad
K_{13}=\frac{r_1+r_3-r_2}{2},\quad
K_{23}=\frac{r_2+r_3-r_1}{2}.
\]
W szczególności $\operatorname{Ric}=0$ wtedy i tylko wtedy,
gdy pełny tensor krzywizny Riemanna jest zerowy.
\end{twierdzenie}
\begin{proof}
Przypomnijmy, że
$r_1=K_{12}+K_{13}$,
$r_2=K_{12}+K_{23}$ i
$r_3=K_{13}+K_{23}$.
Rozwiązanie tego układu liniowego daje podane wzory.
W wymiarze trzy wszystkie składowe tensora krzywizny można
odzyskać z tensora Ricciego (równoważnie: część Weyla jest
identycznie zerowa). W szczególności zerowanie się wszystkich
wartości własnych tensora Ricciego daje zerowanie się tensora
Riemanna.
\end{proof}

\begin{wniosek}
Jeżeli trójwymiarowa rozmaitość Riemanna spełnia
$\operatorname{Ric}=\lambda g$ dla pewnej stałej $\lambda$,
to ma stałą krzywiznę przekrojową $K=\lambda/2$.
Każda zupełna jednospójna trójrozmaitość spełniająca
$\operatorname{Ric}=0$ jest izometryczna z $\mathbb R^3$.
\end{wniosek}
\begin{proof}
Pierwsza teza wynika bezpośrednio z powyższych wzorów.
W drugim przypadku metryka jest płaska; zastosowanie
klasyfikacji zupełnych jednospójnych przestrzeni o stałej
krzywiźnie daje $\mathbb R^3$.
\end{proof}

\begin{uwaga}[Co zmienia się w wyższych wymiarach?]
Dla $n\ge4$ tensor Weyla może być niezerowy, więc metryka
Ricci-płaska nie musi być płaska. Własność wykorzystana
powyżej jest specyficzna dla wymiaru trzy i wyjaśnia,
dlaczego warunek $\operatorname{Ric}=0$ ma w geometrii
trójwymiarowej znacznie silniejsze konsekwencje niż w
geometrii czterowymiarowej.
\end{uwaga}

\section{Hiperboliczne trójrozmaitości i wielościany fundamentalne}
\label{sec:thurston-hyperbolic}

Konstrukcja zwartych ilorazów $\mathbb H^3$ wymaga innej techniki
niż torus odwzorowania czy kratowe ilorazy grup Liego.
Naturalnym narzędziem są wielościany fundamentalne i izometrie
parujące ich ściany.

\begin{definicja}[Dziedzina fundamentalna]
Niech $\Gamma$ będzie dyskretną grupą izometrii $X$.
Domknięty zbiór $P\subset X$ nazywamy \emph{dziedziną
fundamentalną}, jeżeli $\Gamma P=X$, a wnętrza jego różnych
translat są rozłączne, tzn.
$\operatorname{int}(P)\cap\gamma\operatorname{int}(P)=\varnothing$
dla $\gamma\in\Gamma\setminus\{e\}$.
\end{definicja}

Twierdzenie Poincarégo o wielościanach daje warunki, pod którymi
przyporządkowanie izometrii parom ścian wielościanu definiuje
dyskretną grupę, a sam wielościan jest fundamentalny. W pełnej
wersji należy zweryfikować zgodność parowania ścian,
relacje wokół krawędzi i warunki lokalnego domknięcia oraz
zupełności. Sam fakt, że ściany dają się sparować, nie
wystarcza do otrzymania poprawnej rozmaitości hiperbolicznej.

\begin{przyklad}[Dopełnienie węzła ósemkowego]
Niech $K\subset\mathbb S^3$ będzie węzłem ósemkowym.
Wnętrze dopełnienia $\mathbb S^3\setminus N(K)$, gdzie $N(K)$
jest regularnym otoczeniem węzła, ma zupełną strukturę
hiperboliczną o skończonej objętości. Można ją zbudować przez
sklejenie dwóch regularnych idealnych czworościanów
hiperbolicznych. Jeśli $v_3$ oznacza objętość regularnego
idealnego czworościanu, to objętość tej rozmaitości wynosi
$2v_3$. Nie jest to przykład \emph{zwartej} rozmaitości:
na końcu znajduje się cusp (koniec paraboliczny) z
przekrojem będącym dwuwymiarowym torusem.
\end{przyklad}

\begin{przyklad}[Rozmaitość Giesekinga]
Rozmaitość Giesekinga można otrzymać przez odpowiednie
identyfikacje ścian jednego regularnego idealnego
czworościanu hiperbolicznego. Jest ona nieorientowalna,
niezwarta, zupełna i ma skończoną objętość $v_3$.
Jej orientowalne nakrycie dwukrotne jest izometryczne z
dopełnieniem węzła ósemkowego (z naturalną metryką
hiperboliczną). Ten przykład pokazuje, że topologia końców
i orientowalność są istotne nawet przy stałej krzywiźnie $-1$.
\end{przyklad}

\section{Klasyfikacja lokalnych geometrii a geometryzacja}

Nie należy utożsamiać dwóch twierdzeń: klasyfikacja Thurstona
opisuje osiem możliwych \emph{maksymalnych geometrii modelowych}
na jednospójnych jednorodnych przestrzeniach, natomiast
\emph{geometryzacja} dotyczy globalnej struktury topologicznej
dowolnych trójrozmaitości.

\begin{twierdzenie}[Geometryzacja Thurstona--Perelmana, schemat]
Każda zwarta orientowalna trójrozmaitość po rozkładzie na
składniki pierwsze (cięcia wzdłuż sfer dwuwymiarowych),
a następnie odpowiednim rozkładzie wzdłuż nieściśliwych
torusów, daje części typu Seiferta i części hiperboliczne.
Hiperboliczne wnętrza odpowiednich części mają zupełne
metryki hiperboliczne o skończonej objętości.
Wśród możliwych geometrii składników zamkniętych i
przykładów podstawowych pojawia się osiem geometrii
z twierdzenia~\ref{thm:osiem-geometrii}.
\end{twierdzenie}

Sformułowanie jest celowo schematyczne: poprawna wersja
wymaga uwzględnienia redukcji sferycznej, rozkładu JSJ,
przypadków o pustym rozkładzie i sposobu traktowania brzegu.
Dowód geometryzacji nie jest konsekwencją samej listy ośmiu
geometrii. Historycznie przełomem było zastosowanie przepływu
Ricciego z chirurgią; schemat rozwoju metryki zapisuje
równanie
\[
 \frac{\partial g(t)}{\partial t}=-2\operatorname{Ric}(g(t)).
\]
Jego szczegółowa analiza wymaga oddzielnego rozdziału.

\begin{uwaga}[Różne znaczenia słowa „geometria”]
Nie każda zwarta trójrozmaitość dopuszcza jedną geometrię
jednorodną na całej swojej przestrzeni. Geometryzacja mówi,
że po odpowiednich cięciach otrzymujemy części modelowane
przez geometrie z listy. Jest to zasadnicza różnica między
rozmaitością \emph{geometryczną} a rozmaitością
\emph{geometryzowalną}.
\end{uwaga}

\section{Podsumowanie i powiązania z poprzednimi rozdziałami}

Osiem geometrii Thurstona powstaje z dwóch współdziałających
idei: jednorodności metrycznej i istnienia zwartych ilorazów.
Z punktu widzenia rachunku różniczkowego najważniejsze są
globalne ortonormalne ramy, wzór Koszula i obliczenia krzywizn
Nil oraz Sol. Z punktu widzenia topologii zasadniczą rolę
odgrywają grupy podstawowe, nakrycia, wiązki okręgów, liczba
Eulera i torusy odwzorowania. Geometria
$\widetilde{\mathrm{SL}}_2(\mathbb R)$ łączy oba punkty widzenia
przez krzywiznę połączenia wiązki Seiferta.

Warto przy tym pamiętać, że twierdzenie o geometryzacji
ma charakter globalny i nie wynika z lokalnej listy geometrii.
Jest ono pomostem do bardziej zaawansowanej topologii
trójrozmaitości, teorii rozkładów JSJ i przepływu Ricciego.

\paragraph{Literatura uzupełniająca.}
W.~P.~Thurston, \emph{Three-Dimensional Geometry and Topology},
vol.~1, Princeton University Press, 1997;
P.~Scott, \emph{The Geometries of 3-Manifolds},
Bulletin of the London Mathematical Society 15 (1983), 401--487;
J.~W.~Milnor, \emph{Curvatures of Left Invariant Metrics on Lie Groups},
Advances in Mathematics 21 (1976), 293--329;
S.~A.~Pinillos Prado, \emph{Geometric Structures on 3-Dimensional
Manifolds}, Universidade Federal de Minas Gerais, 2021.
